\documentclass[10pt,a4paper]{article}
\usepackage{amsmath,amssymb,amsthm,graphicx}
\usepackage{mathrsfs,bbm,thmtools,thm-restate,xifthen,hyperref,enumitem,mathtools,bm}
\newtheorem{theorem}{Theorem}[section]
\newtheorem{definition}[theorem]{Definition}
\newtheorem{lemma}[theorem]{Lemma}
\newtheorem{claim}[theorem]{Claim}
\newtheorem{question}[theorem]{Question}
\newtheorem{proposition}[theorem]{Proposition}
\newtheorem{fact}[theorem]{Fact}
\newtheorem{corollary}[theorem]{Corollary}
\newtheorem{observation}[theorem]{Observation}
\newtheorem{notation}[theorem]{Notation}
\newtheorem{example}[theorem]{Example}
\newtheorem{remark}[theorem]{Remark}
\newtheorem{note}[theorem]{Note}
\newtheorem*{notation*}{Notation}
\newtheorem*{lemma*}{Lemma}
\newtheorem*{proposition*}{Proposition}
\newtheorem*{note*}{Note}
\newtheorem*{theorem*}{Theorem}
\newtheorem*{corollary*}{Corollary}




% Logic

\newcommand{\innerifequals}[3]{\ifthenelse{\equal{#1}{#2}}{#3}{}}
\newenvironment{switch}[1]{\newcommand{\case}{\innerifequals{#1}}}{}
\newcommand{\definedas}[0]{\coloneqq}

% Linear Algebra
\DeclareMathOperator*{\im}{Im}
\newcommand{\vectorfield}[0]{V}
% \renewcommand{\vec}[1]{\overrightarrow{#1}}


\newcommand{\alignparens}[2]{\left( #1 \vphantom{#2} \right. & \left. \vphantom{#1} #2 \right)}
\newcommand{\alignset}[2]{\left\{ #1 \vphantom{#2} \suchthat \right. & \left. \vphantom{#1} #2 \right\}}
\newcommand{\multilinenorm}[2]{\left\| #1 \vphantom{#2} \right. \\ & \quad \left. \vphantom{#1} #2 \right\|}

\newcommand{\parens}[1]{\left( #1 \right)}
\newcommand{\sparens}[1]{\left[ #1 \right]}
\newcommand{\leftsparens}[1]{\left[ #1 \right.}
\newcommand{\rightsparens}[1]{\left. #1 \right]}
\newcommand{\innerprod}[1]{\left< #1 \right>}

\newcommand{\rhalfopen}[2]{\ensuremath{[#1,#2)}}
\newcommand{\lhalfopen}[2]{\ensuremath{(#1,#2]}}


% Probability
\newcommand{\expectation}[2]{\mathbbm{E}_{#1} \sparens{#2}}
% Set theory
\newcommand{\set}[1]{\left\{ #1 \right\}}
\newcommand{\suchthat}[0]{\middle|}

% Proof parameters
\newcommand{\basefield}[0]{\mathbbm{F}}
\newcommand{\blocklength}[0]{n}
\newcommand{\field}[0]{\basefield^{\blocklength}}
\newcommand{\titlefield}[0]{\texorpdfstring{$\field$}{Fn}}
\newcommand{\degree}[0]{d}
\newcommand{\basefieldsize}[0]{p}

\newcommand{\genfinitefield}[0]{\mathbbm{F}_{q}}

% Constants and indexes
\newcommand{\series}[3]{#1_{#2},\ldots,#1_{#3}}
\newcommand{\seriesdep}[4]{#1_{#3}(#2),\ldots,#1_{#4}(#2)}

% bias
\newcommand{\charfunc}[1]{e \sparens{#1}}
\newcommand{\genbias}[0]{\mu}

\newcommand{\genprime}[0]{p}
% Sets
\NewDocumentCommand{\genset}{o}{\IfNoValueTF{#1}{A}{
    \begin{switch}{#1}
        \case{1}{A}
        \case{2}{B}
        \case{3}{C}
    \end{switch}
}}
% unctions
\NewDocumentCommand{\onvarfunc}{o}{\IfNoValueTF{#1}{f}{
    \begin{switch}{#1}
        \case{1}{f}
        \case{2}{g}
        \case{3}{h}
    \end{switch}
}}
\NewDocumentCommand{\onvarfuncset}{o}{\IfNoValueTF{#1}{\mathfrak{f}}{
    \begin{switch}{#1}
        \case{1}{\mathfrak{f}}
        \case{2}{\mathfrak{g}}
        \case{3}{\mathfrak{h}}
    \end{switch}
}}

\NewDocumentCommand{\genfunc}{o}{\IfNoValueTF{#1}{F}{
    \begin{switch}{#1}
        \case{1}{F}
        \case{2}{G}
        \case{3}{H}
    \end{switch}
}}

\NewDocumentCommand{\genfuncset}{o}{\IfNoValueTF{#1}{\mathfrak{F}}{
    \begin{switch}{#1}
        \case{1}{\mathfrak{F}}
        \case{2}{\mathfrak{G}}
        \case{3}{\mathfrak{H}}
    \end{switch}
}}

\newcommand{\funcdef}[3]{#1:#2\rightarrow#3}
\newcommand{\restrictfunc}[2]{#1|_{#2}}

% Polynomials
\NewDocumentCommand{\genpoly}{o}{\IfNoValueTF{#1}{P}{
    \begin{switch}{#1}
        \case{1}{P}
        \case{2}{Q}
        \case{3}{H}
    \end{switch}
}}
\NewDocumentCommand{\genpolyset}{o}{\IfNoValueTF{#1}{\mathcal{P}}{
    \begin{switch}{#1}
        \case{1}{\mathcal{P}}
        \case{2}{\mathcal{Q}}
        \case{3}{\mathcal{H}}
    \end{switch}
}}

\newcommand{\factor}[0]{\mathcal{B}}
\newcommand{\semrefine}[0]{\succeq_{\text{sem}}}
\newcommand{\semrefineex}[1]{\succeq_{\text{sem}|#1}}
\newcommand{\relsemrefine}[1]{\succeq^{#1}_{\text{sem}}}
\newcommand{\relsemrefineex}[1]{\succeq^{#1}_{\text{sem}}}
\newcommand{\synrefine}[0]{\succeq_{\text{syn}}}
\newcommand{\cubesfactor}[3]{C_{#1,#2}({#3})}

\newcommand{\allpolyset}[3]{Poly_{#1}(#2\rightarrow#3)}

\newcommand{\cube}[2]{(#1|#2)}
\newcommand{\almostcube}[2]{(#1|#2)^{\prime}}
\newcommand{\cubes}[2]{C_{#1}(#2)}
\newcommand{\almostcubes}[3]{C^{\prime}_{{#1},{#2}}({#3})}


\NewDocumentCommand{\lowdegpoly}{o}{\IfNoValueTF{#1}{\alpha}{
    \begin{switch}{#1}
        \case{1}{\alpha}
        \case{2}{\beta}
    \end{switch}
}}

\newcommand{\depth}[1]{depth(#1)}


% Constant Sets
\newcommand{\torus}[0]{\mathbbm{T}}
\newcommand{\naturalnumbersset}[0]{\mathbbm{N}}
\newcommand{\realnumbersset}[0]{\mathbbm{R}}
\newcommand{\complextnumberset}[0]{\mathbbm{C}}


% Varieties
\NewDocumentCommand{\varpoly}{o}{\IfNoValueTF{#1}{L}{
    \begin{switch}{#1}
        \case{1}{L}
        \case{2}{M}
    \end{switch}
}}
\NewDocumentCommand{\varpolyset}{o}{\IfNoValueTF{#1}{\tilde{\mathcal{L}}}{
    \begin{switch}{#1}
        \case{1}{\matchcal{L}}
        \case{2}{\matchcal{M}}
    \end{switch}
}}

\newcommand{\relativeremainder}[1]{\overline{#1}}
\newcommand{\remainderpoly}[0]{R}

\NewDocumentCommand{\onvarpoly}{o}{\IfNoValueTF{#1}{p}{
    \begin{switch}{#1}
        \case{1}{p}
        \case{2}{q}
        \case{3}{h}
    \end{switch}
}}
\NewDocumentCommand{\onvarpolyset}{o}{\IfNoValueTF{#1}{\mathfrak{p}}{
    \begin{switch}{#1}
        \case{1}{\mathfrak{p}}
        \case{2}{\mathfrak{q}}
        \case{3}{\mathfrak{h}}
    \end{switch}
}}

% Specific variety notations
\newcommand{\variety}{{\tilde{X}}}
\newcommand{\titlevariety}{\texorpdfstring{$\variety$}{X}}
%\newcommand{\varietypolycount}{\tilde{c}}
\newcommand{\varietypolycount}{{\tilde{c}}}
\newcommand{\varietydeg}[0]{\tilde{\degree}}
\newcommand{\varietycodeminsion}[0]{\tilde{L}}
\newcommand{\varietyminepsilon}[1]{\tilde{\epsilon}_{#1}}


\newcommand{\lift}[1]{\widehat{#1}}
\newcommand{\homopart}[1]{\mathrm{h}({#1})}
\newcommand{\derivative}[1]{{#1}^{\prime}}
\newcommand{\directionderivative}[0]{D}
\newcommand{\derivativepoly}[1]{\vartheta{#1}}

\newcommand{\zerofunc}[1]{Z \parens{#1}}
\newcommand{\existfunc}[1]{1_{#1}}

% Factors


% Properties
\newcommand{\rank}[1]{rank \parens{#1}}
\newcommand{\drank}[2]{rank_{#1} \parens{#2}}
\newcommand{\rankex}[2]{rank_{#1} \parens{#2}}
\newcommand{\schmrank}[1]{schmrank \parens{#1}}
\newcommand{\lemmarank}[1]{rank^{#1}}
\newcommand{\relschmrank}[2]{shcmrank_{#1}\parens{{#2}}}
\newcommand{\relrank}[2]{rank_{#1} \parens{#2}}
\newcommand{\drelrank}[3]{rank_{{#1}, {#2}} \parens{#3}}


%Constants Notations
\newcommand{\rankval}[0]{{{r}}}
\newcommand{\rankfunc}[0]{{{{r}}}}
\newcommand{\rankbiasfunc}[0]{{\tilde{r}}}
\newcommand{\varietyrankval}[0]{{\bar{r}}}
\newcommand{\dominatedrankfunc}[0]{{r_{\square}}}
\newcommand{\epsilonlimitedrankbias}[0]{{\tilde{\epsilon}}}
\newcommand{\constlimiteddominated}[0]{{c_{\square}}}


\newcommand{\bias}[1]{bias({#1})}
\newcommand{\relbias}[2]{bias_{#1}{({#2})}}

% Norms
\newcommand{\norm}[1]{\left\lVert{#1}\right\rVert}
\newcommand{\abs}[1]{\left| #1 \right|}
\newcommand{\gowersnorm}[2]{{\norm{#2}}_{U^{#1}}}
\newcommand{\infnorm}[1]{{\norm{#1}}_{\infty}}
\newcommand{\dist}[1]{dist \parens{#1}}


%Probability
\newcommand{\pr}[1]{\Pr\left[ #1 \right]}
\newcommand{\prex}[2]{\Pr_{#1}\left[ #2 \right]}
\newcommand\given[1][]{\:#1\vert\:}

% Constraints
\newcommand{\linearformset}[0]{L}
\newcommand{\linearformsetvaribables}[0]{\tilde{\constraintsset}}
\NewDocumentCommand{\linearform}{o}{\IfNoValueTF{#1}{l}{
    \begin{switch}{#1}
        \case{0}{l}
        \case{1}{r}
    \end{switch}
}}
\newcommand{\uniformityerror}[0]{\tau}
\newcommand{\uniformitcomplexity}[0]{\xi}
\newcommand{\density}[0]{\delta}


\newcommand{\gencode}{\mathfrak{C}}
\newcommand{\generatormatrix}{\boldsymbol{G}}
\newcommand{\paritymatrix}{\boldsymbol{H}}

\newcommand{\quotientcode}{\gencode_{\variety}}

\newcommand{\reedmullercode}[3]{RM_{#1}(#2, #3)}
\newcommand{\reedmullercodeex}[3]{RM_{#1, #2}(#3)}
\newcommand{\listpolycount}[4]{\ell_{#1, #2}(#3, #4)}
\newcommand{\listdecodingradiusex}[3]{LDR_{#1, #2}(#3)}
\newcommand{\normalizedcodedistance}[2]{\delta_{#1}(#2)}
\newcommand{\normalizedcodedistanceex}[3]{\delta_{#1, #2}(#3)}
\newcommand{\trycommand}[3][2]{(#2 + #3)^#1}

\DeclareRobustCommand{\topbot}{\genfrac{}{}{0pt}{}}

\title{Distance inheritance for quotient Reed-Muller codes}
\author{Omri Gotlib, Tali Kaufman, Shachar Lovett}
\date{doi:10.4230/LIPIcs.CCC.2025.1; CC BY4.0}
\begin{document}
\maketitle
\setcounter{section}{1}
\noindent\textit{Editorial note.} Original authors: Omri Gotlib, Tali Kaufman and Shachar Lovett. arXiv:2502.15650v1 (2025-02-21), https://arxiv.org/abs/2502.15650v1. Creative Commons Attribution4.0: https://creativecommons.org/licenses/by/4.0/.\par\medskip
\noindent\textit{Editorial note.} Selected complete native author theorem, proof, full supporting sections and bibliography in standalone layout. TeX input files flattened as explicit spans. Unused frontmatter/packages omitted. No mathematical corrections or new proofs.\par\medskip
\noindent\textit{Allowed standard prerequisites (editorial).}\begin{itemize}
\item Finite-field Reed-Muller distance and Bhowmick-Lovett low-complexity approximation/degree preservation. Full author preliminaries and exact cited standard guarantees included.
\end{itemize}\medskip
\section[Preliminaries]{Preliminaries}\label{sec:preliminaries}

\subsection{Basic Definitions and Notations}\label{subsec:definitions_and_notations}
We denote by $\naturalnumbersset$ the set set of integers, i.e natural numbers (excluding 0).
For an integer $k$ we denote $\sparens{k} \definedas \set{1,2,...,k}$.
We use $y = x \pm \epsilon$ to denote $y \in \sparens{x - \epsilon, x + \epsilon}$, and similarly $y = x \mp \lambda$ to denote $y \in \sparens{x + \lambda, x - \lambda}$ (usually when $\lambda < 0$).
\newline
Fix a prime field $\basefield = \basefield_{\genprime}$.
Denote by $\abs{\cdot}$ the natural map of $\basefield$ to $\set{1,...,\genprime - 1} \in \naturalnumbersset$.
We denote the character from $\basefield$ by $\charfunc{x} \definedas e^{2 \pi i \cdot \abs{x}}$

Generally speaking and unless stated otherwise, we use the following conventions:
We use $\blocklength \in \naturalnumbersset$ to denote the number of variables in Reed-Muller code.
We use $\degree$ to denote a degree (typically the degree of the polynomials in our code), and $\variety$ to denote the subset of $\field$ we work in i.e $\variety \subseteq \field$.
Properties of the subset $\variety$ will usually be denoted with $\tilde{\square}$.
We use $\genfunc[1], \genfunc[2], \genfunc[3]$ to denote general functions with domain $\field$, and $\onvarfunc[1], \onvarfunc[2], \onvarfunc[3]$ to denote functions with domain $\variety$.
We use $\genfuncset[1], \genfuncset[2], \genfuncset[3]$ and $\onvarfuncset[1], \onvarfuncset[2], \onvarfuncset[3]$ respectively to denote sets of such functions.
Similarly, we use $\genpoly[1], \genpoly[2], \genpoly[3]$ to denote polynomials with domain $\field$, and $\onvarpoly[1], \onvarpoly[2], \onvarpoly[3]$ polynomials with domain $\variety$.
We use $\genpolyset[1], \genpolyset[2], \genpolyset[3]$ and $\onvarpolyset[1], \onvarpolyset[2], \onvarpolyset[3]$ respectively to denote sets of such polynomials.

\subsection[Polynomials in \titlefield]{Polynomials in \titlefield}\label{subsec:polynomials_in_Fn}
We start by presenting a standard definition for a polynomial over a finite field.
\begin{definition}[Polynomial: Global Definition]
    Let $\degree \in \naturalnumbersset$ be a constant.
    A function $\funcdef{\genpoly}{\field}{\basefield}$ is called a \emph{polynomial of degree $\leq \degree$} if
    it is of the following form:
    \[
        \genpoly(x_1,...,x_{\blocklength}) =
            \sum_{0 \leq \degree_1,...,\degree_{n} : \sum_{i=1}^{\blocklength}{\degree_i} \leq \degree}
                {c_{\degree_1,...,\degree_n} \prod_{i=1}^{\blocklength} x_i^{\degree_i}}
    \]
    We denote the set of all polynomials of degree $\leq \degree$ by $\allpolyset{\leq \degree}{\field}{\basefield}$.
    The value $\degree$ in the definition above is called the \emph{global degree} of the function $\genpoly$, shorthand by the \emph{degree} of $\genpoly$,
    and it is denoted by $\deg(\genpoly) = \degree$.
    \newline
    Additionally, the set of all polynomials from $\field$ to $\basefield$ of degree $\leq \degree$ is denoted by:
    \[
        \allpolyset{\leq \degree}{\field}{\basefield}
    \]
\end{definition}
Next, we present a known equivalent definition for a polynomial using derivatives.
To do so, we first define a derivative in the case of finite fields.
\begin{definition}[Derivative]
    Given a function $\funcdef{\genfunc}{\field}{\basefield}$ and $a \in \field$,
    we define the derivative of $\genfunc$ in direction $a$ as a function $\funcdef{\directionderivative_{a}\genfunc}{\field}{\basefield}$
    defined as follows:
    \[
        \directionderivative_a \genfunc(x) \definedas \genfunc(x + a) - \genfunc(x)
    \]
\end{definition}
\begin{lemma}
    Let $\degree \in \naturalnumbersset$.
    A function $\funcdef{\genfunc}{\field}{\basefield}$ is a polynomial of degree $\leq \degree$
    if and only if $\directionderivative_{a}\genfunc$ is a polynomial of degree $\leq \degree - 1$ for all $a \in \field$.
\end{lemma}
This leads us to a natural definition of a \emph{degree} of a function using derivatives.
\begin{definition}[Local Degree]
    For a function $\funcdef{\genfunc}{\field}{\basefield}$, we define its \emph{local degree},
    to be the least integer $\degree \in \naturalnumbersset$ such that
    for all $a_1,...,a_{\degree + 1}, x \in \field$:
    \[
        \directionderivative_{a_{\degree + 1}}...\directionderivative_{a_{1}} \genfunc(x) = 0
    \]
\end{definition}
In $\field$, the two definitions of degree coincide, and we get a single definition of a degree:
\begin{lemma}[Equivalance of definitions of a degree]\label{lemma:alternative-definition-for-polynomial-using-derivatives}
    Let $\funcdef{\genfunc}{\field}{\basefield}$ be a function, and let $\degree \in \naturalnumbersset$ be an integer.
    Then, the global degree of a $\genfunc$ \emph{equals} its local degree.
\end{lemma}
\begin{remark}
    We sometimes refer to the requirement that the local degree of a function is $\leq \degree$,
    as the \emph{local criteria} of degree $\leq \degree$ polynomials.
\end{remark}

\subsection[Rank-bias in \titlefield]{Rank-bias in \titlefield}
We start by defining the notion of \emph{bias}, which is a measure of how the function is far from being equidistributed (see Appendix~\ref{sec:equidistribution-of-functions} for the exact details).
\begin{definition}[Bias]
    Let $\funcdef{\genfunc}{\field}{\basefield}$.
    The \emph{bias} of the function $\genfunc$ is defined in the following way:
    \[
        \bias{\genfunc} \definedas 1 / \abs{\field} \cdot \sum_{x \in \field}{\charfunc{\genfunc(x)}}
    \]
    Moreover, for a subset $\variety \subseteq \field$, we define the \emph{bias of $\genfunc$ in $\variety$} to be:
    \[
        \relbias{\variety}{\genfunc} \definedas 1 / \abs{\variety} \cdot \sum_{x \in \variety}{\charfunc{\genfunc(x)}}
    \]
\end{definition}
Next, we present a standard definition of rank of a polynomial, which is a notion that measures how \emph{structured} is the function.
Note that low rank implies the polynomial is highly structured.
Formally we have the following definition:
\begin{definition}[Rank of a Polynomial]\label{definition:rank}
    Given a constant $\degree \in \naturalnumbersset$ and a polynomial $\genpoly$,
    the \emph{$\degree$-rank} of $\genpoly$, denoted as $\drank{\degree}{\genpoly}$ is defined to be
    the smallest integer $\rankval$ such that $\genpoly$ can be computed given $\rankval$ polynomials of degree $< \degree$.
    In other wards, we say $\drank{\degree}{\genpoly} = \rankval$ if $\rankval$ is the smallest integer such that
    there exists $\rankval$ polynomials $\genpoly[2]_1,...,\genpoly[2]_\rankval \in \allpolyset{\leq \degree - 1}{\field}{\basefield}$
    and a function $\funcdef{\Gamma}{\field}{\basefield}$ such that:
    \[
        \genpoly(x) = \Gamma \parens{\genpoly[2]_1(x),...,\genpoly[2]_\rankval(x)}
    \]
    If $\degree = 1$, then $1$-rank is defined to be $\infty$ if $\genpoly$ is non constant, and $0$ otherwise.
    \newline
    Moreover, for a polynomial $\genpoly$ of degree $\deg(\genpoly) = \degree$ we denote $\rank{\genpoly} \definedas \drank{\degree}{\genpoly}$.
    \newline
    We call such function $\Gamma$ a \emph{decomposition} or a \emph{computation} of $\genpoly$ using lower-degree polynomials.
\end{definition}
Let us now define a factor.
Note that we focus our discussion to factors in $\field$,
but define the basic definitions over a general set $U$ so they will apply for factors over a general sets.
This is necessary as we will later use them also for other sets such as $\variety \subseteq \field$.
\begin{definition}[Factor]
    Let $U$ be a set.
    A \emph{factor over $U$}, denoted by $\factor$, is simply a partition of the set $U$.
    Each subset in the partition is called an \emph{atom}.
    \newline
    A collection of function $\funcdef{\genfunc_1,...,\genfunc_c}{U}{\basefield}$ defines a factor $\factor_{\genfunc_1,...,\genfunc_c}$ over $U$
    with atoms:
    \[
        \set{u \in U \suchthat f_1(u)=b_1, ..., f_c(u)=b_c}
    \]
    for all $b_1,...,b_c \in \basefield$.
    \newline
    Additionally, we use $\factor$ to also denote the map $\factor(u) \rightarrow (\genfunc_1(u),...,\genfunc_c(u))$.
\end{definition}

\begin{notation*}
    Let $\funcdef{\genfunc_1,...,\genfunc_c}{U}{\basefield}$ be a collection of functions.
    For a factor $\factor \definedas \factor_{\genfunc_1,...,\genfunc_c}$, we denote by $\abs{\factor}$ the amount of functions that define it, i.e. $\abs{\factor} = c$.
    Moreover, we denote $\norm{\factor} \definedas \abs{\basefield}^c$, which is the maximal amount of (possibly empty) atoms.
    Additionally, the rank of the factor is defined to be the rank of the polynomials that generate it.
\end{notation*}
\begin{definition}[Polynomial Factor]
    We say a factor $\factor$ over $\field$ is a \emph{polynomial factor} if it is defined by a collection of polynomials $\funcdef{\genpoly_1,...,\genpoly_c}{\field}{\basefield}$,
    i.e. $\factor = \factor_{\genpoly_1,...,\genpoly_c}$.
    The degree of the factor, denote as $\deg(\factor)$ is the maximal degree of the polynomials $\genpoly_1,...,\genpoly_c$.
\end{definition}
Note that the notion of degree (and polynomial) are defined only for functions over $\field$,
therefore this definition is well-defined only for $U = \field$.

%TODO: Avoid using "factor" before it is defined.
\begin{definition}[Rank of a Factor]\label{definition:factor-rank}
    Let $\genpolyset$ be a collection of polynomials $\funcdef{\genpoly_1,...,\genpoly_c}{\field}{\basefield}$.
    The rank of the polynomial collection is defined as:
    \[
        \rank{\genpolyset} \definedas \min \set{\drank{\degree}{\sum_{i=1}^c{\lambda_i \genpoly_i}} \suchthat 0 \neq \vec{\lambda} \in \basefield^c, \degree = \max_{i\in\sparens{c}}{\deg(\lambda_i \genpoly_i)}}
    \]
    For a factor $\factor$ defined by a collection of polynomials $\genpolyset$, we define its rank to be the rank of the collection of polynomials defining it.
    For a non-decreasing function $\funcdef{\rankfunc}{\naturalnumbersset}{\naturalnumbersset}$, a factor $\factor$ is called $\rankfunc$-regular if its rank is at least $\rankfunc(\abs{\factor})$.
\end{definition}
\begin{note*}
    Note that in the definition above, the rank of each linear combination is calculated as the $\degree$-rank,
    where $\degree$ is the maximal degree of a polynomial that participates in the linear combination non-trivially.
    This is crucial as it ensures that a high rank factor do not have linear dependence in the largest-degree homogenous component of any of its polynomials.
\end{note*}


We now present a fundamental property of high rank polynomials, that was first proved by~\cite{green2007distribution} when $\degree < \abs{\basefield}$,
later extended to general fields by~\cite{kaufman2008worst}, and further extended also to large fields by~\cite{DBLP:journals/corr/0001L15}.
This property of high rank polynomials is that they they have low bias:
\begin{theorem}[Rank-bias in $\field$]\label{high-rank-implies-low-bias}
    Let $\basefield$ be a finite field.
    Let $\epsilon > 0$ and $\degree \in \naturalnumbersset$.
    There exists $r_{\ref{high-rank-implies-low-bias}} \definedas r_{\ref{high-rank-implies-low-bias}}(\basefield, \degree, \epsilon)$,
    such that for every degree-$\degree$ polynomial $\funcdef{\genpoly}{\field}{\basefield}$:
    if $\rank{\genpoly} \geq r_{\ref{high-rank-implies-low-bias}}$ then $\bias{\genpoly} < \epsilon$.
\end{theorem}
\begin{remark}
    This property implies that a collection of polynomials that have high rank is \emph{equidistributed}.
    See Appendix~\ref{sec:equidistribution-of-functions} for more details in this regard.
\end{remark}

%
%Moreover, high rank polynomial factors do not just have low bias, they also have low $\degree$-gowers norm:
%\begin{theorem}\cite[Theorem 1.20]{tao2011inverse}\label{high-rank-implies-low-gowers-norm}
%    Let $\basefield$ be a finite field.
%    Let $\epsilon > 0$ and $\degree \in \naturalnumbersset$.
%    There exists $r_{\ref{high-rank-implies-low-gowers-norm}} \definedas r_{\ref{high-rank-implies-low-gowers-norm}}(\basefield, \degree, \epsilon)$ such that for every non-classical polynomial $\funcdef{\genpoly}{\field}{\basefield}$
%    of degree $\leq \degree$, if $\rank{\genpoly} \geq r_{\ref{high-rank-implies-low-gowers-norm}}$ then $\gowersnorm{\degree}{\charfunc{\genpoly}} < \epsilon$
%\end{theorem}
%
%Furthermore, when looking on the gowers norm, the other direction is also true: uniform polynomials have high rank:
%\begin{theorem}\cite[Claim 7.16]{book}\label{low-gowers-norm-implies-high-rank}
%    Let $\degree, r \geq 1$ be integers.
%    There exists $\epsilon_{\ref{low-gowers-norm-implies-high-rank}}(\basefield, \degree, r)$ such that every non-classical polynomial $\funcdef{\genpoly}{\field}{\basefield}$
%    of degree $\leq \degree$, if $\gowersnorm{\degree}{\charfunc{\genpoly}} < \epsilon_{\ref{low-gowers-norm-implies-high-rank}}$ then $\rank{\genpoly} \geq r$
%\end{theorem}
%
%%TODO: Delete torus and write to F.


\subsection[Regularization in \titlefield]{Regularization in \titlefield}\label{subsec:regularization-in-Fn}
In this subsection we define the regularization process in $\field$.
Before doing so, let us present some definitions in this regard.
Note that we define the basic definitions over a general set $U$ so they will apply for factors over a general sets,
as this is necessary as we will later use them also for other sets such as $\variety \subseteq \field$.
\begin{definition}[Measureable]
    Let $U$ be a set, and let $A \subseteq U$.
    Let $\genfuncset = \set{\genfunc_1,...,\genfunc_c}$ be a collection of functions $\funcdef{\genfunc_i}{U}{\basefield}$.
    We say a function $\funcdef{\genfunc[2]}{U}{\basefield}$ is \emph{measurable in respect of $\genfuncset$ in $A$}, shorthand by \emph{$\genfuncset$-measurable in $A$},
    if there exists a function $\funcdef{\Gamma}{\basefield^c}{\basefield}$ such that:
    \[
        \forall a \in A: g(a) = \Gamma(\genfunc_1(a),...,\genfunc_c(a))
    \]
    When discussing the factor over $A$ defined by $\factor = \factor_{\genfunc_1,...,\genfunc_c}$,
    we also say $\genfunc[2]$ is \emph{measurable in resepct of $\factor$}.
    The function $\Gamma$ will be denoted as the \emph{measurement function} of $\genfunc[2]$ in respect of $\genfuncset$.
    Additionally, when $A = U$, we sometimes omit the specification of the domain, and say $\genfunc[2]$ is measurable in respect of $\genfuncset$.
    \newline
    Note that in this paper, we usually think of $U = \field$, and $A$ is either $\field$ or $\variety \subseteq \field$.
\end{definition}
\begin{remark}
    If $\genfunc[2]$ is $\genfuncset$-measurable in $A$, then every value of $\genfunc[2]$ in $A$
    can be determined by the values of $\genfunc_1,...,\genfunc_c$.
    In other words, the function $\genfunc[2]$ is constant inside every atom of the factor defined by $\genfuncset$.
\end{remark}

\begin{definition}[Syntactic Refinement]
    Let $\factor$ and $\factor^\prime$ be polynomial factors over $U$.
    We say a factor $\factor^\prime$ is a \emph{syntactic refinement} of the factor $\factor$, if the collection of functions defining $\factor$ is a subset of the set
    of functions defining $\factor^\prime$.
    We denote this property of $\factor^\prime$ by $\factor^\prime \synrefine \factor$.
\end{definition}
We now present a standard generalized definition of refinement, where we only require the atoms induced by the refined factors are sub-atoms of those that are induced by the original factor.
Note that in this refinement, we allow the refined factor to include completely different polynomials than the original factor.
\begin{definition}[Semantic Refinement]\label{def:semantic-refinement}
    Let $\factor$ and $\factor^\prime$ be polynomial factors on $U$ defined by $\genpolyset$ and $\genpolyset^{\prime}$ respectively.
    We say the factor $\factor^\prime$ is a \emph{semantic refinement} of the factor $\factor$ in $A \subseteq U$,
    if $x, y \in A$ with $\factor^\prime(x) = \factor^\prime(y)$ implies that $\factor(x)=\factor(y)$.
    We denote this property of $\factor^\prime$ by $\factor^\prime \semrefineex{A} \factor$.
    When $A = U$, we sometimes omit $A$ from the syntax and denote it with $\factor^{\prime} \semrefine \factor$
    \newline
    Note that $\factor^\prime \synrefine \factor$ implies $\factor^\prime \semrefineex{A} \factor$ for every $A \subseteq U$.
\end{definition}
\begin{remark}
    A handy property of semantic refinement is that if $\funcdef{\genfunc}{A}{\basefield}$ is $\genpolyset$-measurable,
    then it is also $\genpolyset^{\prime}$-measurable in $A$.
    Moreover, the other direction is also true:
    If every $\genpolyset$-measurable function $\funcdef{\genfunc}{A}{\basefield}$ in $A$ is also $\genpolyset^{\prime}$-measurable in $A$, then $\factor^{\prime} \semrefineex{A} \factor$.
\end{remark}

Next, we recall a lemma that was presented in~\cite[Theorem 4.1]{bhattacharyya2013locally}, that allows us, given a polynomial that is measurable by a high rank factor in $\field$,
to replace the polynomials in the measurement function to any collection of polynomials with smaller or equal degree, and preserve the degree of the original polynomial.
Note that we state the lemma under the constraint that $\degree < \abs{\basefield}$, but it is also valid for when $\degree \geq \abs{\basefield}$ with proper generalization of definitions to this case (See~\cite[Theorem 4.1]{bhattacharyya2013locally} for the exact statement).
\begin{lemma}[Preserving Degree in $\field$]\label{preserving-degree-starting-field}
Let $\degree>0$ an integer such that $\degree < \abs{\basefield}$, and let $\funcdef{\genpoly_1,...,\genpoly_c}{\field}{\basefield}$ be polynomials of degree at most $\degree$, that form a factor of rank $\geq  \rankfunc^{\ref{preserving-degree-starting-field}}(\basefield, \degree, c)$.
Assume that for $\funcdef{\Gamma}{\basefield^{c}}{\basefield}$, the function $\funcdef{\gamma}{\field}{\basefield}$ defined as $\gamma(a) \definedas \Gamma(\genpoly_1(a),...,\genpoly_c(a))$ is of $\deg(\gamma)=\degree^{\prime}$.
\newline
Then, for every collection of polynomials $\funcdef{\genpoly[2]_1,...,\genpoly[2]_c}{\field}{\basefield}$ that satisfy $\deg(\genpoly[2]_i) \leq \deg(\genpoly_{i})$,
the function $\gamma^{\prime}$ defined as $\gamma^{\prime}(a)=\Gamma(\genpoly[2]_1(a),...,\genpoly[2]_c(a))$ is a polynomial of $\deg(\gamma^{\prime})\leq\degree^{\prime}$.
\end{lemma}
Next, we restate a useful lemma from~\cite[Lemma 4.17]{DBLP:journals/corr/0001L15} that shows that under the conditions
above, $\Gamma$ is as a low-degree polynomial (with even stronger conditions).
Formally, they showed:
\begin{lemma}[Faithful Composition]
    In the case discussed above, the structure of $\Gamma$ is as follows:
    \[
        \Gamma(z_1,...,z_{c_1})  =
        \sum_{\alpha \in \sparens{\basefieldsize - 1}^{c}} {{C_{\alpha}} \cdot {\prod_{i = 1}^{c}}{z_i^{\alpha_i}}}
    \]
    where $C_{\alpha} = 0$ whenever $\sum_{i = 1}^{c_1}(\alpha_i \cdot \deg(\genpoly_i)) > \degree^\prime$.
    \newline
    In other words, this means that $\Gamma$ as a function $\funcdef{\Gamma}{\basefield^{c}}{\basefield}$,
    is a polynomial of degree $\leq \degree^{\prime}$, even when substituting its $i$-th input by any polynomial of degree $\leq \deg(\genpoly_i)$.
\end{lemma}


Finally, we restate the regularization process, that was first presented by~\cite[Lemma 2.3]{green2007distribution}.
The regularization process shows that every factor have a high-rank factor that semantically refines it,
without increasing the size of the factor too much (its new size is independent of $\blocklength$).
\begin{lemma}[Regularization in $\field$]\label{regularization-in-Fn-lemma}
    Let $\funcdef{\rankfunc}{\naturalnumbersset}{\naturalnumbersset}$ be a non-decreasing function and let $\degree \in \naturalnumbersset$.
    There exists $\funcdef{C_{\rankfunc, \degree}^{\ref{regularization-in-Fn-lemma}}}{\naturalnumbersset}{\naturalnumbersset}$ such that the following holds:
    Let $\factor$ be a factor on $\field$ defined by polynomials $\genpolyset = (\genpoly_1,...,\genpoly_c)$ where for all $i \in [c]$: $\funcdef{\genpoly_i}{\field}{\basefield}$ and $\deg(\genpoly_i) \leq \degree$.
    Then, there is an $\rankfunc$-regular factor $\factor^\prime$ defined by polynomials $\genpolyset[2] = (\genpoly[2]_1,...,\genpoly[2]_{c^\prime})$ where
    for all $i \in [c]$: $\funcdef{\genpoly[2]_i}{\field}{\basefield}$ and $\deg(\genpoly[2]_i) \leq \degree$ such that
    $\factor^\prime \semrefine \factor$ and $c^\prime \leq C_{\rankfunc, \degree}^{\ref{regularization-in-Fn-lemma}}(c)$.
    \newline
    Moreover, if $\factor \synrefine \bar{\factor}$ for some polynomial factor $\bar{\factor}$ with rank at least $\rankfunc(c^\prime)+c^\prime+1$,
    then we can require that $\factor^\prime \synrefine \bar{\factor}$.
\end{lemma}



\section[Polynomials in \titlevariety]{Polynomials in \titlevariety}\label{sec:polynomials_in_X}
In this section we wish to generalize the definition of degree-$\degree$ polynomials for functions $\funcdef{\onvarfunc}{\variety}{\basefield}$.
Note that we wish to define it using a property of $\onvarfunc$ that is intrinsic to $\variety$: given a function $\funcdef{\onvarfunc}{\variety}{\basefield}$,
we wish be able to determine its degree only using values of $\variety$, without considering any value outside of $\variety$ (such as values of $\field \setminus \variety$).
\newline
To define such property, we generalize the local definition of a degree that is defined for polynomials in $\field$.
We remind the reader that in $\field$, we said a function over $\field$ is a polynomial of degree $\leq \degree$ if and only if its $(\degree+1)$-derivative in every direction is $\equiv 0$.
Thus, in order to determine the $(\degree+1)$-derivative of a function in directions $y_1,...,y_{\degree+1}$, one needs to evaluate the function over all the points of the cube generated by $x,y_1,...,y_{\degree+1}$,
which is the set of points $\set{x + \sum_{i \in S}{y_i}}_{S \subseteq [\degree+1]}$.
This raises a challnge in extending this definition for functions defined over $\variety \subseteq \field$:
depending on $\variety$, the function $\funcdef{\onvarfunc}{\variety}{\basefield}$ is not be defined to all points in all the cubes of $\field$, because some of those points do not lie in $\variety$.
\newline
Therefore, to generalize the definition of a polynomial to $\variety$, we start by giving the formal definition and notation of the set of cubes in $\variety$:
\begin{definition}[Cubes]
    Let $k \in \naturalnumbersset$ be an integer and let $x, y_1,...,y_{k} \in \field$.
    We define the cube $\cube{x}{y_1,...y_{k}}$ as follows:
    \[
        \cube{x}{y_1,...,y_{k}} \definedas \set{x + \sum_{i \in S}{y_i}}_{S \subseteq [k]}
    \]
    We refer to $x$ as the \emph{offset} of the cube, and $y_1,...,y_{k}$ as the \emph{directions} of the cube.
    \newline
    Moreover, Let $\variety \subseteq \field$ be a subset.
    We define the \emph{set of cubes of $\variety$ of size $k$} as follows:
    \[
        \cubes{k}{\variety} \definedas
            \set{\cube{x}{y_1,...,y_{k}} \suchthat \forall S \subseteq[k]: (x + \sum_{i \in S}{y_i}) \in \variety}
    \]
\end{definition}

Using this definition, we can define a polynomial of degree $\leq \degree$ for subsets of $\field$:
\begin{definition}[Polynomials in $\variety$]
    Let $\degree \in \naturalnumbersset$ be an integer, and let $\variety \subseteq \field$.
    We say the degree of a function $\funcdef{\onvarfunc}{\variety}{\basefield}$ is $\degree$
    if $\degree$ is the smallest integer such that $\onvarfunc$ vanishes over all cubes of size $(\degree+1)$, i.e:
    \[
        \forall {\cube{x}{y_1,...,y_{\degree + 1}} \in \cubes{\degree+1}{\variety}}:
            \directionderivative_{y_{\degree + 1}}...\directionderivative_{y_1} \onvarpoly(x) = 0
    \]
    A function over $\variety$ of degree $\leq \degree$ is also called a \emph{polynomial of degree $\leq \degree$}.
    We denote the set of polynomials of degree $\leq \degree$ over $\variety$ by $\allpolyset{\leq \degree}{\variety}{\basefield}$.
\end{definition}
\begin{note*}
    For $\variety = \field$, the definition above coincides with the local definition of polynomials.
\end{note*}

\subsection[Lifting Polynomials]{Lifting Polynomials}\label{subsec:lifting-polynomials}
Our goal to achieve good properties for polynomials over $\variety$.
To do so, we wish to connect the desired properties of polynomials defined over $\variety$, to properties known for polynomials over $\field$.
Following such strategy raises a question: given a polynomial $\funcdef{\onvarpoly}{\variety}{\basefield}$, which polynomial over $\field$ should we consider to deduce properties of $\onvarpoly$?
To find such a polynomial over $\field$, it would have been useful that all polynomials over $\variety$ actually "came from" polynomials over $\field$.
More formally, it would have been useful that all polynomials $\funcdef{\onvarpoly}{\variety}{\basefield}$ would be equal to a restriction of some polynomial $\funcdef{\genpoly}{\field}{\basefield}$ of degree $\leq \degree$, to the set $\variety$.
This would give us a "good candidate" (or candidates) to polynomials over $\field$, that using their known properties, we could achieve the properties we desire for polynomials over $\variety$.
\newline
Generally speaking, the existence of such polynomial $\funcdef{\genpoly}{\field}{\basefield}$ is not trivial by itself, and it mapy depend on the polynomial $\onvarpoly$ and the set $\variety$.
In this subsection, we discuss sets $\variety \subseteq \field$ that have this property for every polynomial $\funcdef{\onvarpoly}{\variety}{\basefield}$.
Before formulating the notion above, we start by a simple remark:
\begin{remark}
    By the local criteria for $\field$, we have that a restriction of a polynomial of degree $\leq \degree$ over $\field$ to $\variety$ is a polynomial of degree $\leq \degree $ over $\variety$.
    Therefore, the other direction is true: every restriction of a polynomial over $\field$ to $\variety$ is a polynomial over $\variety$.
\end{remark}
%First, let us present an exmaple that shows that not every set has this property~\cite[Example 1.4]{kazhdan2019extendingweaklypolynomialfunctions}:
%\begin{example}
%    Fix $\blocklength = 2$.
%    Define $\varpoly(x_1, x_2) = x_1 \cdot x_2 \cdot (x_1 - x_2)$,
%    and consider:
%    \[
%        \variety \definedas \zerofunc{\varpoly} = \set{x = (x_1, x_2) \in \basefield^2 \suchthat x_1 = 0 \vee x_2 = 0 \vee x_1 = x_2}
%    \]
%    Now, consider the function $\funcdef{\onvarpoly}{\variety}{\basefield}$ defined as:
%    $\onvarpoly(x_1, 0) = \onvarpoly(0, x_2) = 0, \onvarpoly(x_1, x_1) = x_1$.
%
%\end{example}

Next, let us define subsets $\variety \subseteq \field$ that have the desired property, which we call \emph{$\degree$-lift-enabler variety}.
\begin{definition}[$\degree$-lift-enabler Subset]
    Let $\basefield$ be a field, and $\blocklength>0$ be an integer.
    For an integer $\degree > 0$, we say a subset $\variety\subseteq\field$ is \emph{$\degree$-lift-enabler} if for every $\degree^{\prime} \leq \degree$,
    for every polynomial $\onvarpoly \in \allpolyset{\degree^{\prime}}{\variety}{\basefield}$
    there exist a polynomial $\lift{\onvarpoly} \in \allpolyset{\degree^{\prime}}{\field}{\basefield}$ such that $\restrictfunc{\onvarpoly}{\variety}=\restrictfunc{\lift{\onvarpoly}}{\variety}$.
\end{definition}
\begin{remark}
    Using the local criterion of polynomials and the fact that that $\cubes{\degree+1}{\variety} \subseteq \cubes{\degree+1}{\field}$,
    one can see that for a polynomial $\funcdef{\onvarpoly}{\variety}{\basefield}$ with $\deg(\onvarpoly) = \degree$,
    every extension $\funcdef{\genpoly}{\field}{\basefield}$ with $\onvarpoly = \restrictfunc{\genpoly}{\variety}$
    holds the bound $\deg(\lift{\onvarpoly}) \geq \degree$.
    The other direction is not true in the general case, but it is specifically promised when the variety is $\degree$-lift-enabler.
\end{remark}

This definition naturally raises the following definition:
\begin{definition}[The Lift Operator]
    Let $\degree \in \naturalnumbersset$ be an integer.
    Let $\variety \subseteq \field$ be a $\degree$-lift-enabler subset.
    We define \emph{the $\degree$-lift operator} to be an operator $\funcdef{\lift{\square}}{\allpolyset{\leq \degree}{\variety}{\basefield}}{\allpolyset{\leq \degree}{\field}{\basefield}}$ the following way:
    \newline
    Let $\degree^\prime \leq \degree$.
    Given a polynomial $\funcdef{\onvarpoly}{\variety}{\basefield}$ of degree $\degree^\prime$, the operator $\lift{\square}$ returns a polynomial $\funcdef{\lift{\onvarpoly}}{\field}{\basefield}$ of degree $\degree^\prime$
    such that $\onvarpoly = \restrictfunc{\lift{\onvarpoly}}{\variety}$.
    Note that we did not require the lift to be unique.
    Thus, in case there are multiple valid lifts for a polynomial $\onvarpoly \in \allpolyset{\leq \degree}{\variety}{\basefield}$, the lift operator picks a single (consistent) one of them.
    Moreover, the lift always exists because the subset $\variety$ is $\degree$-lift-enabler.
    \newline
    In addition, for a collection $\onvarpolyset = (\onvarpoly_1,...,\onvarpoly_c)$ of polynomials $\onvarpoly_i \in \allpolyset{\leq \degree}{\variety}{\basefield}$,
    we denote $\lift{\onvarpolyset} \definedas (\lift{\onvarpoly_1},...,\lift{\onvarpoly_c})$
\end{definition}

In the following subsections, we give example to two concrete sets $\variety \subseteq \field$ that are $\degree$-lift-enablers.
Before doing so, we define an algebraic variety:
\begin{definition}[Algebraic Variety]
    For a collection of functions $\genfuncset \definedas \set{\genfunc_1,...\genfunc_c}$ such that $\funcdef{\genfunc_i}{\field}{\basefield}$,
    we denote $\zerofunc{\genfuncset} \definedas \set{x \in \field \suchthat \forall i: \genfunc_i(x) = 0}$.
    \newline
    If the collection is a collection of polynomials, we call $\zerofunc{\genfuncset}$ an \emph{algebraic variety}, shorthand by \emph{variety}.
    \newline
    The degree of the variety is defined to be the maximal degree of polynomials in the collection that defines it.
    The rank of the variety is defined to be the rank of the collection that defines the variety (as a collection).
\end{definition}

\subsection[High Rank Varieties of High Minimal Degree]{High Rank Varieities of High Minimal Degree}\label{subsec:high-rank-varieities-of-high-degree}
We now present a theorem proved in \cite[Corollary 1.10]{kazhdan2018polynomial}, that shows that high rank varieties are $\degree$-lift-enabler when the polynomials defining the variety are of degree $>\degree$:
\begin{theorem}\label{subsec:high-rank-varities-are-d-lift-enabler}
    Let $\basefield$ be a finite field, and let $\varietydeg$, $\varietypolycount>0$ representing parameters of a variety.
    Let $\degree < \varietydeg$ a positive integer representing a degree of a polynomial which we wish to lift.
    There exists $\varietyrankval=\varietyrankval(\basefield,\varietydeg,\varietypolycount)>0$ such that for for all $\blocklength \in \naturalnumbersset$, any variety $\variety = \zerofunc{\varpolyset} \subseteq \field$ for $\varpolyset = (\varpoly_1,...,\varpoly_{\varietypolycount})$
    such that $\rank{\varpolyset} > \varietyrankval$, degree $\deg(\varpolyset) = \varietydeg$, with all defining polynomials of degree $\deg(\varpoly_i)> \degree$,
    it holds that $\variety$ is a $\degree$-lift-enabler subset.
\end{theorem}
\begin{remark}
    Under the conditions stated above, it was proved in ~\cite{kazhdan2018polynomial} that the lift is in fact \emph{unique}.
    Formally, if $\funcdef{\onvarpoly}{\variety}{\basefield}$ is a polynomial of degree $\leq \degree$,
    then there exists a \emph{unique} polynomial $\funcdef{\genpoly }{\field}{\basefield}$ such that $\restrictfunc{\genpoly}{\variety} \equiv \onvarpoly$.
\end{remark}

\subsection[High Rank Varieties on a Large Field]{High Rank Varieties on a Large Field}\label{subsec:high-rank-varieties-on-a-large-field}
In this subsection, we recall a theorem proved by~\cite[Theorem 1.7]{kazhdan2019extendingweaklypolynomialfunctions} regarding high rank varieties that are defined on "large" fields.
We note that the fields are large in respect of the degree $\degree$ one wish to lift, but still does not depend on $\blocklength$.

Next we define a weakly polynomial, which generalizes our definition of a polynomial in $\variety$, that was used in~\cite[Definition 1.1]{kazhdan2019extendingweaklypolynomialfunctions}:
\begin{definition}
    Let $\variety \subseteq \field$ be a set.
    We say a function $\funcdef{\genfunc}{\variety}{\basefield}$ is a \emph{weakly polynomial of degree $\leq \degree$}
    if for any affine subspace $L \subseteq \variety$, the restriction $\restrictfunc{\genfunc}{L}$ is a polynomial of degree $\leq \degree$.
\end{definition}
\begin{remark}
    By the local criteria of a polynomial, it is easy to see that every $\genpoly \in \allpolyset{\leq \degree}{\variety}{\basefield}$ is a weakly polynomial of degree $\leq \degree$.
\end{remark}
And now, we can present the lifting theorem for large fields, as proved in~\cite[Theorem 2.17]{kazhdan2020propertieshighranksubvarieties}.
\begin{theorem}\label{high-rank-varieties-over-large-fields-are-d-lift-enablers}\cite[Theorem 2.17]{kazhdan2020propertieshighranksubvarieties}
    Let $\degree, \varietydeg \in \naturalnumbersset$,
    and let $\basefield$ be a finite field such that $\abs{\basefield} > \degree \cdot \varietydeg$.
    There exists $\rankfunc_{\ref{high-rank-varieties-over-large-fields-are-d-lift-enablers}} = \rankfunc_{\ref{high-rank-varieties-over-large-fields-are-d-lift-enablers}}(\varietydeg, \degree)$
    such that for any variety $\variety \subseteq \field$ of maximal degree $\leq \varietydeg$ and rank $\geq \rankfunc_{\ref{high-rank-varieties-over-large-fields-are-d-lift-enablers}}$,
    have the following property:
    Every weakly polynomial function $\funcdef{\onvarpoly }{\variety}{\basefield}$ of degree $\leq \degree$
    can be lifted to a polynomial function $\funcdef{\genpoly}{\field}{\basefield}$ of degree $\leq \degree$.
\end{theorem}
\begin{note*}
    Note that we stated the theorem above to finite fields, but it is also valid for infinite algebraically closed fields.
\end{note*}
The theorem above implies the following corollary:
\begin{corollary}
    Let $\degree, \varietydeg \in \naturalnumbersset$,
    and let $\basefield$ be a finite field such that $\abs{\basefield} > \degree \cdot \varietydeg$.
    There exists $\rankfunc_{\ref{high-rank-varieties-over-large-fields-are-d-lift-enablers}} = \rankfunc_{\ref{high-rank-varieties-over-large-fields-are-d-lift-enablers}}(\varietydeg, \degree)$
    such that for any variety $\variety \subseteq \field$ of degree $\varietydeg$ and rank $\geq \rankfunc_{\ref{high-rank-varieties-over-large-fields-are-d-lift-enablers}}$
    is a $\degree$-lift-enabler.
\end{corollary}


\section[Relative Rank-Bias Property]{Relative Rank-Bias Property}\label{sec:relative-rank-bias-property}
In this section, we generalize the relation between rank and bias that is known for $\field$ also for $\variety \subseteq \field$.
Specifically, in Lemma~\ref{high-rank-implies-low-bias}, it was shown that high rank factors have low bias in $\field$.
We wish to define an alternative definition of rank for $\variety \subseteq \field$, called \emph{$\variety$-relative rank},
such that high $\variety$-relative rank implies low bias in $\variety$.
This type of relation (and definition) was shown previously to a few sets;
in~\cite[Theorem 1.8]{lampert2021relative} for sets $\variety = \zerofunc{\genpolyset[2]}$ where $\genpolyset[2]$ is a collection of polynomials of high rank;
and in~\cite[Theorem 1.4]{gowers2022equidistributionhighrankpolynomialsvariables} for sets $\variety = S^\blocklength$ for $S \subset \basefield$.

To understand this notion, we first introduce a simple example that demonstrates the need for a different definition of rank to achieve equidistribution properties in subsets of $\field$.
\begin{example}
    Let $\variety = \set{x \in \field \suchthat x_1 = 0}$.
    Define $\funcdef{\genpoly}{\field}{\basefield}$ by $\genpoly(x) \definedas x_1$.
    \newline
    In $\field$, $\genpoly$ has rank $\infty$ as it can not be decomposed polynomials of degree $< 1$ (constants).
    Additionally, it is perfectly equidistributed.
    This is the simplest example of the rank-bias relation in $\field$.
    \newline
    However, when restricting $\genpoly$ to $\variety$, we get $\restrictfunc{\genpoly}{\variety} \equiv 0$.
    As $0$ is a constant function, it is the \emph{least} equidistributed possible in $\variety$.
    Therefore, we see that the way we defined rank in $\field$ does not imply the desired equidistribution in $\variety$:
    we found a polynomial with high rank (infinity) that has a very high bias \emph{in $\variety$} (the maximal).
\end{example}
The reason the definition of rank in $\field$ fails to capture equidistribution even on subsets that are really similar to $\field$ (isomorphic to $\basefield^k$), is because of the following reason:
Even though our polynomial $\genpoly$ does \emph{not} have a decomposition to a few lower-degree polynomials by itself,
there \emph{exists} a \emph{$\variety$-equivalent} polynomial that has such structured decomposition.
Here, by $\variety$-equivalent we mean a polynomial in $\field$ that is bounded by the same degree bound, and is equal to $\genpoly$ in $\variety$.
In the example described above, this equivalent polynomial is the constant function $0$, and its decomposition is the trivial one (any function decomposes a constant function).
An alternative perspective which we use throughout this paper to $\variety$-equivalence is that both polynomials are equal up to a \emph{valid $\variety$-remainder}: a bounded degree polynomial that is $\equiv 0$ in $\variety$.

Generally speaking, high $\variety$-relative rank may not imply low bias in $\variety$.
Therefore, this structure-randomness relation is not true for a general subset $\variety \subseteq \field$,
but is a \emph{property} of the subset $\variety$.
Thus, we say that a subset has the \emph{relative rank-bias property} if this relation holds,
i.e. if high $\variety$-relative rank implies equidistribution in $\variety$.

Let us now formally define our definition for relative rank, inspired by the two different definitions of relative rank presented in~\cite[Definition 1.6]{lampert2021relative}
and in~\cite[Definition 1.3]{gowers2022equidistributionhighrankpolynomialsvariables}:
\begin{definition}[Relative Rank of a Polynomial]\label{def:relative-rank-of-polynomial}
    Let $\variety \subseteq \field$ and let $\degree \in \naturalnumbersset$.
    For an integer $\degree \in \naturalnumbersset$ and a polynomial $\funcdef{\genpoly}{\field}{\basefield}$, we define its \emph{$\degree$-relative rank} in respect of $\variety$ as:
    \[
        \drelrank{\degree}{\variety}{\genpoly} \definedas \min \set{\drank{\degree}{\genpoly - \relativeremainder{\genpoly}} \suchthat
        \relativeremainder{\genpoly} \in \allpolyset{\leq \deg(\genpoly)}{\field}{\basefield}, \restrictfunc{\relativeremainder{\genpoly}}{\variety} \equiv 0}
    \]
    For a polynomial $\genpoly$ of degree $\deg(\genpoly) = \degree$ we denote $\relrank{\variety}{\genpoly} \definedas \drelrank{\degree}{\variety}{\genpoly}$.
\end{definition}
\begin{definition}[$\variety$-equivalent and $\variety$-remainder]
    Moreover, we say a polynomial is \emph{$\variety$-equivalent} to $\genpoly$ if its restriction to $\variety$ is $\equiv \restrictfunc{\genpoly}{\variety}$.
    We say it is \emph{valid $\variety$-equivalent to $\genpoly$} if it is $\variety$-equivalent to $\genpoly$ and its of the \emph{same} degree of $\genpoly$.
    \newline
    Similarly, we say a polynomial is \emph{$\variety$-remainder} of $\genpoly$ if its restriction to $\variety$ is \emph{$\equiv 0$}.
    We say it is \emph{valid $\variety$-remainder} if it is $\variety$-remainder of $\genpoly$ and its of degree \emph{smaller or equal} of the degree of $\genpoly$.
    \newline
    %TODO: Denote also the equivalent polynomial and make the notations in the other sections consistent with it.
    We typically denote such polynomial as $\relativeremainder{\genpoly}$.
\end{definition}
In other words, the $\degree$-relative rank of a polynomial $\genpoly$ is the smallest $\degree$-rank of all valid $\variety$-equivalents of $\genpoly$.

\begin{note}\label{note:comparison-to-gowers-rank}
    Note that~\cite[Definition 1.3]{gowers2022equidistributionhighrankpolynomialsvariables} defines rank in a substantially different way than our definition,
    and consequentially our results will not apply to the sets they presented.
    One of the main differences in the definition of rank occurs for $\degree = 1$.
    In the definition we use for rank, the rank of every (non-constant) degree-$1$ polynomial is $\infty$, where in the definition used in~\cite{gowers2022equidistributionhighrankpolynomialsvariables} it is a finite number (which is possibly very small).
    This difference is crucial, as for example, it makes regularization according to their definition not-trivially possible,
    where it is known to be possible when rank is defined by the definition we use (Lemma~\ref{regularization-in-Fn-lemma}).
\end{note}
\begin{definition}[Relative Rank of a Factor]
    Let $\variety \subseteq \field$.
    Let $\genpolyset$ be a set of polynomials $\genpolyset = \set{\genpoly_1,...,\genpoly_c}$.
    The rank of the polynomial set $\genpolyset$ relative to the subset $\variety$ is defined as:
    \[
        \relrank{\variety}{\genpolyset} \definedas \min \set{\drelrank{\degree}{\variety}{\sum_{i=1}^c{\lambda_i \genpoly_i}} \suchthat 0 \neq \vec{\lambda} \in \basefield^c, \degree = \max_{i \in \sparens{c}}{\deg(\lambda_i \genpoly_i)}}
    \]
    For a factor $\factor$ defined by a collection of polynomials, we define its relative rank relative to $\variety$ to be the relative rank of the collection of polynomials defining it, relative to the set $\variety$.
    \newline
    For a non-decreasing function $\funcdef{\rankfunc}{\naturalnumbersset}{\naturalnumbersset}$, a factor $\factor$ is called $\rankfunc$-$\variety$-regular if its relative rank in respect to $\variety$ is at least $\rankfunc(\abs{\factor})$.
\end{definition}

\subsection[Relative Rank-Bias Property]{Relative Rank-Bias Property}
%TODO: Add introduction to this subsection
\begin{definition}[Relative Rank-Bias property]
    Let $\basefield$ be a finite field, and let $\degree \in \naturalnumbersset$ be an integer.
    Let $\funcdef{\rankbiasfunc}{\realnumbersset^{+}}{\naturalnumbersset}$ be a function that represents the rank-bias relation for a fixed $\degree, \basefield$.
    \newline
    We say a set $\variety \subseteq \field$ has the \emph{$(\rankbiasfunc, \basefield, \degree)$-relative rank-bias property} if
    for every $\epsilon > 0$,
    for every polynomial $\genpoly$ of degree $\leq \degree$ with $\relrank{\variety}{\genpoly} \geq \rankbiasfunc(\epsilon)$ we have:
    \[
        \relbias{\variety}{\genpoly} < \epsilon
    \]
\end{definition}

As an immediate corollary of Lemma~\ref{high-rank-implies-low-bias} that shows that high rank implies low bias, we have that $\variety = \field$ has the relative rank-bias property.
\begin{corollary}[$\field$ has the relative rank-bias property]
    For every finite field $\basefield$ and $\degree \in \naturalnumbersset$, let $\funcdef{\rankbiasfunc}{\realnumbersset^{+}}{\naturalnumbersset}$ defined as $\rankbiasfunc(\epsilon) \definedas \rankval_{\ref{high-rank-implies-low-bias}}(\basefield, \degree, \epsilon)$.
    Then, we have that the set $\variety = \field$ has the $(\rankbiasfunc, \basefield, \degree)$-relative rank-bias property.
\end{corollary}
\begin{proof}
    This is a simple usage of Lemma~\ref{high-rank-implies-low-bias}:
    Note that when $\variety = \field$, we have that $\relrank{\variety}{\genpoly} = \rank{\genpoly}$.
    Now, if $\genpoly$ is a polynomial of degree $\leq \degree$ and $\relrank{\variety}{\genpoly} = \rank{\genpoly} \geq \rankbiasfunc(\epsilon) = \rankbiasfunc_{\ref{high-rank-implies-low-bias}}(\basefield, \degree, \epsilon)$,
    then:
    \[
        \relbias{x \in \field}{\genpoly(x)} < \epsilon
    \]
\end{proof}


\subsection[Limited Relative Rank-Bias Property]{Limited-Relative Rank-Bias Property}\label{subsec:limited-relative-rank-bias-property}
Sometimes, however, we can not request $\variety$ to be such that high relative rank implies low bias for every $\epsilon > 0$,
but only for $\epsilon^\prime \geq \epsilon$ for some constant $\epsilon > 0$.
This leads to defining the \emph{limited relative rank-bias property}, which will be used to discuss such sets $\variety \subseteq \field$.
\newline
As we will later see, this definition raises naturally where $\variety$ is a high rank variety,
in which for the relative rank-bias property to hold for some $\epsilon > 0$, the rank of the variety should be greater than a value that is dependent of $\epsilon$.
Thus, to have the relative rank-bias property for a high rank variety but without requiring an infinitely large rank, we must limit the relative rank-bias property for $\epsilon^\prime \geq \epsilon$
We formulate the definition of this property as follows:
\begin{definition}[Limited Relative Rank-bias property]
    Let $\basefield$ be a finite field, let $\degree \in \naturalnumbersset$ be an integer, and let $\epsilon > 0$ be a constant.
    Let $\funcdef{\rankbiasfunc}{[\epsilon, \infty]}{\naturalnumbersset}$ be a function that represents the limited-relative-rank-bias relation.
    \newline
    We say a set $\variety \subseteq \field$ has the \emph{$(\rankbiasfunc, \basefield, \degree, \epsilon)$-limited-relative-rank-bias property} if
    for every $\epsilon^\prime \geq \epsilon$,
    for every polynomial $\genpoly$ of degree $\leq \degree$ with $\relrank{\variety}{\genpoly} \geq \rankbiasfunc(\epsilon^\prime)$ we have:
    \[
        \relbias{\variety}{\genpoly} < \epsilon^{\prime}
    \]
    As a convention, we denote by $\epsilonlimitedrankbias$ the $\epsilon$ such that the limited-relative-rank-bias property holds for $\variety$.
\end{definition}

%TODO: Change relative rank to be rank_X instead of rank_L.
\subsubsection[High Rank Varieities]{High Rank Varieities}
In this subsection, we are discussing specifically $\variety \subseteq \field$ that are in the form $\variety = \zerofunc{\genpolyset[2]}$ for a set of polynomials $\genpolyset[2]$ that form a high rank factor.
Let us present some known results of the relative rank-bias relation for high rank varieites:
In the scenario when we are working relative to $\variety$, the equivalent for Theorem~\ref{high-rank-implies-low-bias} is also known when we assume $\degree < char(\basefield)$, as shown in~\cite[Theorem 1.8]{lampert2021relative}:
\begin{theorem}[High relative rank implies low bias in high rank varieties]\label{high-relative-schmidt-rank-implies-low-relative-bias}
    Let $\basefield$ be a finite field and let $0 \leq \degree < char(\basefield)$.
    Let $\epsilon > 0$ be a constant, and let $\varietypolycount \in \naturalnumbersset$.
    There exist $\varietyrankval^{\ref{high-relative-schmidt-rank-implies-low-relative-bias}} = \varietyrankval^{\ref{high-relative-schmidt-rank-implies-low-relative-bias}}(\basefield, \degree, \varietypolycount, \epsilon)$
    and $\rankval^{\ref{high-relative-schmidt-rank-implies-low-relative-bias}} = \rankval^{\ref{high-relative-schmidt-rank-implies-low-relative-bias}}(\basefield, \degree, \epsilon)$ such that the following holds:
    \newline
    Let $\varpolyset = (\varpoly_1,...,\varpoly_{\varietypolycount})$ be a collection of polynomials of degrees $\leq \degree$ with $\rank{\varpolyset} \geq \varietyrankval^{\ref{high-relative-schmidt-rank-implies-low-relative-bias}}$ and let $\genpoly$ be a polynomial of degree $\leq \degree$.
    \newline
    Then, if $\relrank{\varpolyset}{\genpoly} \geq \rankval^{\ref{high-relative-schmidt-rank-implies-low-relative-bias}}$, we have:
    \[
        \relbias{\varpolyset}{\genpoly} < \epsilon
    \]
\end{theorem}
\begin{note*}
    Note that the original statements in~\cite{lampert2021relative} are stated for a different definition of rank, noted as \emph{schmidt rank}.
    In the appendix~\ref{sec:comparing-ranks} we compare the two different definitions, and show that our definition of rank is comprehensive
    enough in a sense that a polynomial with high rank also has high schmidt rank.
    Additionally, we show that for a given $\rankval \in \naturalnumbersset$, the lower bound of rank required for a polynomial to be of schmidt rank $\geq \rankval$,
    is only $c \cdot \rankval$ for some constant $c \in \naturalnumbersset$.
\end{note*}
\begin{remark}
    Note that in the original statement of theorem~\ref{high-relative-schmidt-rank-implies-low-relative-bias} as stated in~\cite[Theorem 1.8]{lampert2021relative},
    there are good bounds on the rank needed for $\varpolyset$ and $\genpoly$ for the theorem to hold.
    \newline
    Specifically, there exist constants $A(\degree), B(\degree)$ such that for an error $\epsilon = \abs{\basefield^{-s}}$,
    if $\varietyrankval^{\ref{high-relative-schmidt-rank-implies-low-relative-bias}} = A(\varietypolycount + s)^B$ and $\rankval^{\ref{high-relative-schmidt-rank-implies-low-relative-bias}} = A(1 + s)^B$,
    then we have:
    \[
        \relbias{\zerofunc{\varpolyset}}{\genpoly} < \abs{\basefield}^{-s}
    \]
    \newline
    In our proof, it is enough that the bounds on $\rankval$ and $\varietyrankval$ are independent of $\blocklength$, thus we omit the exact bounds stated above and use the statement as stated in Theroem~\ref{high-relative-schmidt-rank-implies-low-relative-bias}.
\end{remark}

\begin{remark}\label{in-high-relative-rank-implies-low-relative-bias-epsilon-increasing-rank-requirment-decreasing}
Note that both $\rankval^{\ref{high-relative-schmidt-rank-implies-low-relative-bias}}(\basefield, \degree, \epsilon)$
and $\varietyrankval^{\ref{high-relative-schmidt-rank-implies-low-relative-bias}}(\basefield, \degree, \varietypolycount, \epsilon)$
are decreasing when $\epsilon$ is increasing.
This means for example, that for all $\epsilon^\prime \geq \epsilon$,
a variety that satisfies the theorem's rank condition for $\epsilon$ also satisfies the theorem's rank condition for $\epsilon^\prime$.
Therefore, a polynomial with rank $\geq \rankval^{\ref{high-relative-schmidt-rank-implies-low-relative-bias}}(\basefield, \degree, \epsilon^\prime)$
will have a bias $< \epsilon^\prime$.
\end{remark}
As a corollary of Theorem~\ref{high-relative-schmidt-rank-implies-low-relative-bias} and Corollary~\ref{relative-schimdt-rank-equals-schmidt-rank-if-the-variety-is-of-high-degree}, we have that
high rank varieties has the limited-relative-rank-bias property.
Formally, we have:
\begin{corollary}[High Rank Varieties Have the Limited-Relative Rank-Bias Property]\label{high-rank-variety-has-limited-rank-relative-bias-property}
    Let $\basefield$ be a finite field, and let $\varietydeg \in \naturalnumbersset$ such that $0 < \varietydeg < \abs{\basefield}$.
    Let $\epsilonlimitedrankbias > 0$ be a constant which represents the desired relative rank-bias limit.
    There exists $\funcdef{\rankbiasfunc_{\ref{high-rank-variety-has-limited-rank-relative-bias-property}}}{[\epsilonlimitedrankbias, \infty]}{\naturalnumbersset}$ with $\rankbiasfunc_{\ref{high-rank-variety-has-limited-rank-relative-bias-property}} \definedas \rankbiasfunc_{\ref{high-rank-variety-has-limited-rank-relative-bias-property}}(\basefield, \varietydeg)$
    such that the following holds:
    \newline
    Let $\varietypolycount \in \naturalnumbersset$ be an integer.
    There exists $\varietyrankval_{\ref{high-rank-variety-has-limited-rank-relative-bias-property}} \definedas \varietyrankval_{\ref{high-rank-variety-has-limited-rank-relative-bias-property}}(\basefield, \varietydeg, \varietypolycount, \epsilonlimitedrankbias)$ such that
    for every $\varpolyset = (\varpoly_1,...,\varpoly_{\varietypolycount})$ polynomial factor of degree $\leq \varietydeg$ with $\rank{\varpolyset} \geq \varietyrankval$, we have:
    \newline
    The variety $\variety = \zerofunc{\varpolyset}$ has the $(\rankbiasfunc_{\ref{high-rank-variety-has-limited-rank-relative-bias-property}}, \basefield, \varietydeg, \epsilonlimitedrankbias)$-limited-relative-rank-bias property.
\end{corollary}
\begin{proof}
    Let $\basefield$ be a finite field, and let $\varietydeg \in \naturalnumbersset$ such that $0 < \varietydeg < \abs{\basefield}$.
    Let $\epsilonlimitedrankbias > 0$.
    We choose:
    \[
        \rankbiasfunc_{\ref{high-rank-variety-has-limited-rank-relative-bias-property}}(\epsilon) \definedas
        \rankval_{\ref{high-relative-schmidt-rank-implies-low-relative-bias}}(\basefield, \varietydeg, \epsilon)
    \]
    Note that for every $\epsilon$ in its domain, $\rankbiasfunc_{\ref{high-rank-variety-has-limited-rank-relative-bias-property}}$ does not depend on $\epsilonlimitedrankbias$.
    Let $\varietypolycount \in \naturalnumbersset$.
    Now, we choose:
    \[
        \varietyrankval_{\ref{high-rank-variety-has-limited-rank-relative-bias-property}}(\basefield, \varietydeg, \varietypolycount, \epsilonlimitedrankbias) \definedas
        \varietyrankval_{\ref{high-relative-schmidt-rank-implies-low-relative-bias}}(\basefield, \varietydeg, \varietypolycount, \epsilonlimitedrankbias)
    \]
    Using Theorem~\ref{high-relative-schmidt-rank-implies-low-relative-bias} that shows high rank implies low bias in $\variety$,
    and the assumption that $\epsilon \geq \epsilonlimitedrankbias$ (specifically Remark~\ref{in-high-relative-rank-implies-low-relative-bias-epsilon-increasing-rank-requirment-decreasing})
    concludes the proof.
\end{proof}


\section[Regularization Relative to \titlevariety]{Regularization Relative to \titlevariety}\label{sec:regularization-relative-to-X}

In this section, we generalize the definitions and statements regarding factors and regularization in $\field$,
to their corresponding definitions and statements to relative rank in respect of $\variety \subseteq \field$.
\newline
Note that in oppose to the previous chapter that we discussed a general $U$ and $A \subseteq U$,
in this chapter we discuss only $U = A = \field$.
This is done for clearance and to avoid defining definitions we will not use in our main proof.
\begin{definition}[Measurable Relative to $\variety$]
    Let $\genfuncset = \set{\genfunc_1,...,\genfunc_c}$ be a set of functions $\funcdef{\genfunc_i}{\field}{\basefield}$.
    We say a function $\funcdef{\genfunc[2]}{\field}{\basefield}$ is \emph{measurable in respect of $\genfuncset$ relative to $\variety$},
    or \emph{$\variety$-relative $\genfuncset$-measurable},
    if there exists a function $\funcdef{\relativeremainder{\genfunc[2]}}{\field}{\basefield}$ with $\restrictfunc{\relativeremainder{\genfunc[2]}}{\variety} \equiv 0$
    and a function $\funcdef{\Gamma}{\basefield^c}{\basefield}$ such that:
    \[
        \forall a \in A: \genfunc[2](a) = \Gamma(\genfunc_1(a),...,\genfunc_c(a)) + \relativeremainder{\genfunc[2]}(a)
    \]
    And:
    \[
        \deg(\genfunc[2] - \relativeremainder{\genfunc[2]}) \leq \deg(\genfunc[2])
    \]
    We sometimes refer to $\Gamma$ as the \emph{$\variety$-relative measurement function}.
\end{definition}

\begin{note*}
    Note that if $\deg(\genfunc[2] - \relativeremainder{\genfunc[2]}) \leq \deg(\genfunc[2])$ as discussed above,
    then the same bound  also bounds the degree of the remainder, i.e. $\deg(\relativeremainder{\genfunc[2]}) \leq \deg(\genfunc[2])$.
    Therefore $\relativeremainder{\genfunc[2]}$ is a valid $\variety$-remainder of $\genfunc[2]$.
    Moreover, this requirement is equivalent to the definition above,
    as if $\deg(\relativeremainder{\genfunc[2]}) \leq \deg(\genfunc[2])$, then we also have $\deg(\genfunc[2] - \relativeremainder{\genfunc[2]}) \leq \deg(\genfunc[2])$.
\end{note*}
\begin{note*}
    Also note that without the bound on the degree of the remainder,
    being measurable relative to $\variety$ is in fact equivalent for being a measurable in $A = \variety$.
    This is true because under these conditions, the remainder $\relativeremainder{\genfunc[2]}$ has no constraints but $\restrictfunc{\relativeremainder{\genfunc[2]}}{\variety} \equiv 0$,
    thus the condition left on the measurement is just being a measurement to $\genfunc[2]$ in $\variety$.
\end{note*}
\begin{remark}
    If $\genfunc[2]$ is a function that it is $\genfuncset$-measurable relative to $\variety$,
    then every value of $\genfunc[2]$ can be determined by the values of $\genfunc_1,...,\genfunc_c$ up to a remainder $\relativeremainder{\genfunc[2]}$ of degree $\leq \degree$.
    Thus, perhaps we do not know that the function $\genfunc[2]$ is constant inside every atom of $\genfuncset$ as in a regular semantic refinement,
    but we do know that there exists a function $(\genfunc[2] - \relativeremainder{\genfunc[2]})$ that equals to $\genfunc[2]$ on $\variety$, is constant on every atom of $\genfuncset$ and it is a function with a bounded degree i.e. $\deg(\genfunc[2] - \relativeremainder{\genfunc[2]})\leq \deg(\genfunc[2])$.
\end{remark}

Next, we present a new type of refinement, which is a relaxation of semantic refinement.
This relaxation will allow us to discuss the corresponding claim of the polynomial regularity lemma (Lemma~\ref{regularization-in-Fn-lemma}) for relative rank (instead of rank).
\begin{definition}[Semantic Refinement Relative to $\variety$]
    Let $\factor$ and $\factor^\prime$ be polynomial factors on $\field$, defined by sets of polynomials $\genpolyset, \genpolyset^{\prime}$ respectively,
    and let $\degree \in \naturalnumbersset$.
    We say a factor $\factor^\prime$ is a \emph{semantic refinement relative to $\variety$} of the factor $\factor$,
    or \emph{$\variety$-relative semantic refinement},
    if the following holds:
    Every function $\funcdef{\genfunc}{\field}{\basefield}$ that is $\genpolyset$-measurable relative to $\variety$,
    is also $\genpolyset^{\prime}$-measurable relative to $\variety$.
    If the definition above holds, we denote $\factor^{\prime} \relsemrefineex{\variety} \factor$.
\end{definition}
\begin{note*}
    It is easy to see that this relation is transitive, i.e. if
    $\factor^{\prime} \relsemrefineex{\variety} \factor$ and
    $\factor^{\prime\prime} \relsemrefineex{\variety} \factor^{\prime}$,
    then $\factor^{\prime\prime} \relsemrefineex{\variety} \factor$.
\end{note*}
\begin{remark}
    In $\variety$, semantic refinements relative to $\variety$ behave the same as regular semantic refinements in the perspective of being measurable:
    every function that is $\genpolyset$-measurable in $\variety$ is also $\genpolyset^{\prime}$-measurable in $\variety$.
    However, the two definitions behave differently in the perspective of being measurable in $\field$.
    Specifically, in relative semantic refinements,
    if $\genfunc[2]$ is a $\genpolyset$-measurable function it is not necessarily $\genpolyset^{\prime}$-measurable.
    However, it is measurable up to a remainder $\relativeremainder{\genfunc[2]}$ of degree $\leq \deg(\genfunc[2])$ such that $\restrictfunc{\relativeremainder{\genfunc[2]}}{\variety} \equiv 0$.
\end{remark}
\begin{corollary}\label{relative-semantic-refinement-is-restricted-semantic-refinement}
    If $\factor^{\prime} \relsemrefineex{\variety} \factor$, then in $\variety$ it is a regular semantic refinement, i.e. $\factor^{\prime} \semrefineex{\variety} \factor$.
\end{corollary}

Next, we present a new regularization process that allows us to increase the \emph{relative} rank of a factor without increasing the size of the factor too much (independent of $\blocklength$).
This regularization process generalizes the regularization process in $\field$, which was first presented by~\cite[2.3]{green2007distribution}.
We call this type of regularization process a \emph{relative-regularization process} relative to $\variety$, shorthand by $\variety$-regularization
For a specific function $\rankfunc$, we will sometimes call applying this lemma a \emph{$\rankfunc$-$\variety$-regularization}.
Note that to allow such a relative-regularization process to hold, we must use the relaxed definition of semantic refinement that is presented above.
\begin{theorem}\label{theorem:regularization-in-X}
    Let $\funcdef{\rankfunc}{\naturalnumbersset}{\naturalnumbersset}$ be a non-decreasing function and let $\degree \in \naturalnumbersset$.
    There exists $\funcdef{C_{\rankfunc, \degree}^{\ref{theorem:regularization-in-X}}}{\naturalnumbersset}{\naturalnumbersset}$ such that the following holds:
    Let $\factor$ be a factor defined by polynomials $\genpolyset = (\genpoly_1,...,\genpoly_c)$ where for all $i \in [c]$: $\funcdef{\genpoly_i}{\field}{\basefield}$ and $\deg(\genpoly_i) \leq \degree$.
    Then, there is an $\rankfunc$-$\variety$-regular factor $\factor^\prime$ defined by polynomials $\genpolyset^{\prime} = (\genpoly^{\prime}_{1},...,\genpoly^{\prime}_{c^\prime})$ where
    for all $i \in [c]$: $\funcdef{\genpoly^{\prime}_i}{\field}{\basefield}$ and $\deg(\genpoly^{\prime]}_{i}) \leq \degree$ such that
    $\factor^\prime \relsemrefineex{\variety} \factor$ and $c^\prime \leq C_{\rankfunc, \degree}^{\ref{theorem:regularization-in-X}}(c)$.
    \newline
    Moreover, if $\factor \synrefine \bar{\factor}$ for some polynomial factor $\bar{\factor}$ with
    relative rank of at least $\rankfunc(c^\prime)+c^\prime+1$ and rank of at least ${\rankfunc_{\ref{preserving-degree-starting-field}}(\basefield, \degree, c^{\prime})} + c^\prime + 1$,
    then we can require that $\factor^\prime \synrefine \bar{\factor}$.
\end{theorem}
\begin{proof}
    We follow the lines of the proof given by~\cite{book}[Lemma 7.29], but here, we wish to increase the \emph{relative} rank of the factor instead of its rank.
    We present an iterative process, which will eventually lead us to a factor of size $c^{\prime}$ with relative rank higher than $\rankfunc(c^\prime)$,
    that is a semantic refinement relative to $\variety$.
    Let $\degree \in \naturalnumbersset$,
    and let $\factor$ be a polynomial factor defined by $\genpolyset = (\genpoly_1,...,\genpoly_c)$ such that $\funcdef{\genpoly_i}{\field}{\basefield}$ of degree $\leq \degree$.
    Define $M(\factor) \definedas (M_{\degree},...,M_1) \in \naturalnumbersset^{\degree}$,
    where $M_i$ denotes the number of polynomials in $\genpolyset$ that have degree exactly $i$.
    Thus, $\sum_{i=1}^{\degree}M_i = c$.
    We define the lexicographical order on $\naturalnumbersset^{\degree}$ where $M > M^{\prime}$ if and only if $M_i > M^{\prime}_i$ for some $1 \leq i \leq \degree$,
    and $M_j = M^{\prime}_j$ for all $j > i$.
    This proof will be by transfinite induction on $M$ under the lexicographical order.
    Next we describe a step of the regularization process.
    \newline
    Let $\factor$ be a polynomial factor defined by $\genpolyset = (\genpoly_1,...,\genpoly_c)$.
    Note that this is an abuse of notations: the factor $\factor$ and the set $\genpolyset$ refer to the original factor in the first step, and also to the current factor in the middle of the relative-regularization process.
    If $\factor$ is $\rankfunc$-$\variety$-regular, then we are done.
    Otherwise, we change $\factor$ as follows:
    First, we denote $\rankfunc_{\ref{preserving-degree-starting-field}}^{\basefield, \degree}(c) \definedas \rankfunc_{\ref{preserving-degree-starting-field}}(\basefield, \degree, c)$,
    and we $\rankfunc_{\ref{preserving-degree-starting-field}}$-regularize $\genpolyset$ using lemma~\ref{regularization-in-Fn-lemma}
    to get a set of polynomials $\genpolyset_1 = (\genpoly^1_1,...,\genpoly^1_{c_1})$ of degree $\leq \degree$,
    which defines a factor $\factor_1$ and has a rank $\geq \rankfunc_{\ref{preserving-degree-starting-field}}^{\basefield, \degree}(c_1)$.
    %TODO: Do I need to prove the regular regularization in order to show that?
    Note that $M(\factor_1) \leq M(\factor)$.
    Then, again, if somehow $\factor_1$ is now $\rankfunc$-$\variety$-regular, we are done.
    \newline
    Otherwise, by definition, there exists some linear combination of the polynomials in $\genpolyset_{1}$ that
    has $\degree^\star$-relative rank less than $\rankfunc(c_1)$,
    where $\degree^\star$ is the maximal degree that participates in the linear combination.
    Let $\vec{\genpoly}(x) = \sum_{i=0}^{c_1}{\lambda_i \genpoly^{1}_i(x)}$ where $\vec{0} \neq \vec{\lambda}\in \basefield^{c_1}$,
    be the linear combination with $\drelrank{\degree^\star}{\variety}{\vec{\genpoly}} \leq \rankfunc(c_1)$ where $\degree^\star \definedas \max_{i \in \sparens{c_1}}{\deg(\lambda_i \genpoly^1_i)}$.
    By definition of relative rank, there exists $\relativeremainder{\genpoly} \in \allpolyset{\leq \deg(\vec{\genpoly})}{\field}{\basefield}$ with $\restrictfunc{\relativeremainder{\genpoly}}{\variety} \equiv 0$ such that
    $\drank{\degree^\star}{\vec{\genpoly} - \relativeremainder{\genpoly}} \leq \rankfunc(c_1)$.
    Note that $\deg(\relativeremainder{\genpoly}) \leq \degree^\star$.
    By definition of $\degree^\star$-rank, we have that we can decompose $\vec{\genpoly} - \relativeremainder{\genpoly}$ as a function of $\rankfunc(c_1)$ polynomials of degree $\leq \degree^\star - 1$.
    In other words, there exist a measurement function $\funcdef{\vec{\Gamma}}{\basefield^{\rankfunc(c_1)}}{\basefield}$ and polynomials $\genpoly[2]_1,...,\genpoly[2]_{\rankfunc(c_1)}$
    with $\deg(\genpoly[2]_i) \leq \degree^\star - 1$ such that:
    \[
        \forall a \in \field: \vec{\genpoly}(a) - \relativeremainder{\genpoly}(a) = \vec{\Gamma} \parens {\genpoly[2]_1(a),...,\genpoly[2]_{\rankfunc(c_1)}(a)}
    \]
    Now, let $\genpolyset^{\star} \subseteq \genpolyset_1$ be the set of all such maximal-degree polynomials,
    and let $i^{\star}$ be chosen such that $\genpoly^{1}_{i^\star} \in \genpolyset^{\star}$.
    Note that the set $\genpolyset^\star$ is non empty, as by definition, $\degree^{\star}$ is the maximal degree of polynomial in the expression $\sum_{i=1}^{c_1}{\lambda_{i} \genpoly^{1}_i}$ such that $\lambda_i \neq 0$.
    \newline
    For the next step, define the polynomial factor $\factor_2$ be the polynomial factor defined by the set:
    \[
        \genpolyset_2 \definedas \genpolyset_1 \setminus \set{\genpoly^1_{i^\star}} \cup \set{\genpoly[2]_1,...,\genpoly[2]_{\rankfunc(c_1)}}
    \]
    Finally, the factor $\factor_2$ will be the factor returned from the relative-regularization step.
    \newline
    It is easy to see that if the process above halts, we get a $\rankfunc$-$\variety$-regular factor.
    Now, we prove the first part of the lemma by showing the following claims:
    \begin{claim}
        The factor generated from the regularization above is of bounded size: a bound that may depend on $\rankfunc, \degree, c$, but does not depend on $\blocklength$.
        Formally, we claim that there exists $\funcdef{C^{\ref{theorem:regularization-in-X}}_{\rankfunc, \degree}}{\naturalnumbersset}{\naturalnumbersset}$
        such that we have $c^{\prime} \leq C^{\ref{theorem:regularization-in-X}}_{\rankfunc, \degree}(c)$.
    \end{claim}
    \begin{proof}
        It is enough to prove the following:
        \begin{enumerate}
            \item~\label{relative-regularization-step-factor-size-is-bounded}
            In each step, the amount of polynomials there are in $\genpolyset_1, \genpolyset_2$ are bounded by a bound that depend only on $\rankfunc, \degree, c$ (independent of $\blocklength$).
            \item~\label{relative-regularization-amount-of-steps-is-bounded}
            The number of steps of the relative-regularization process is also bounded by a bound that depends only on $\rankfunc, \degree, c$ (independent of $\blocklength$).
        \end{enumerate}
        The combination of these two will obtain the desired bound of the amount of polynomials in the last-step regularized factor, which is $C^{\ref{theorem:regularization-in-X}}_{\rankfunc, \degree}(c)$.
        Note that the bound on the last-step relative-regularized factor in not simply the multiplication of the two bounds,
        but a recursively-substitution of the bound in~\ref{relative-regularization-step-factor-size-is-bounded},
        a bounded amount of times (bounded by the bound in~\ref{relative-regularization-amount-of-steps-is-bounded}).
        \newline
        For~\ref{relative-regularization-step-factor-size-is-bounded}, we first notice that the number of polynomials in the regular regularization process is bounded,
        specifically we have $\abs{\genpolyset^1} = c_{1} \leq C^{\ref{regularization-in-Fn-lemma}}_{\rankfunc_{\ref{preserving-degree-starting-field}}, \degree}(c)$.
        Moreover, the polynomial factor $\factor_2$ is generated by adding at most $\rankfunc(c_1)$ polynomials to the factor, and thus we have $\abs{\genpolyset_2} \leq c_1 + \rankfunc(c_1)$ which is also bounded by substituting the bound on $c_1$.
        \newline
        For~\ref{relative-regularization-amount-of-steps-is-bounded}, we use the transfinite induction on $M$ we mentioned earlier to show that the process must halt after a bounded number of steps.
        Formally, we show that there exist $M^{\prime}$ which depends only on $M(\factor)$ such that $M(\factor_2) \leq M^{\prime} < M(\factor)$.
        This will bound the number of steps by a value that depend only on $M(\factor)$, which depends only on $\rankfunc, \degree, c$.
        To do so, we first notice that the regular regularization does not increase the value of $M$, i.e. $M(\factor_1) \leq M(\factor)$.
        Thus, we can focus on the second part of the relative-regularization.
        In this part, we replace a single degree $\degree^{\star}$ polynomial by at most $\rankfunc(c_1)$ polynomials of degree $\leq \degree^{\star} - 1$.
        Therefore, by choosing $M^{\prime} \definedas (M_{\degree}, ..., M_{\degree^{\star}+1}, M_{\degree^{\star}}-1, M_{\degree^{\star}-1}+\rankfunc(c_1),...,M_1+\rankfunc(c_1))$
        we get that $M(\factor_2) \leq M^{\prime} < M(\factor_1) \leq M(\factor)$, which concludes~\ref{relative-regularization-amount-of-steps-is-bounded}.
    \end{proof}
    \begin{claim}
        The factor generated from the regularization above is a $\variety$-relative semantic refinement of the original factor, i.e $\factor^\prime \relsemrefineex{\variety} \factor$.
    \end{claim}
    \begin{proof}
        It is enough to show that in each step, the factors generated by the relative-regularization process are semantic refinements relative to $\variety$ of the previous step's factor.
        Specifically, we show $\factor_{2} \relsemrefineex{\variety} \factor_{1} \relsemrefineex{\variety} \factor$ and the claim will follow from transitivity of relative semantic refinements.
        \newline
        We start by proving $\factor_1 \relsemrefineex{\variety} \factor$.
        Let $\funcdef{\genfunc}{\field}{\basefield}$ be a function that is $\genpolyset$-measurable relative to $\variety$.
        We denote $\degree_{\genfunc} \definedas \deg(\genfunc)$.
        By definition, there exists $\funcdef{\Gamma}{\basefield^c}{\basefield}$,
        $\funcdef{\relativeremainder{\genfunc}}{\field}{\basefield}$ where $\deg(\relativeremainder{\genfunc}), \deg(\genfunc - \relativeremainder{\genfunc}) \leq \degree_{\genfunc}$ and $\restrictfunc{\relativeremainder{\genfunc}}{\variety} \equiv 0$, such that:
        \[
            \forall a \in \field: \genfunc(a) = \Gamma(\genpoly_1(a),...,\genpoly_c(a)) + \relativeremainder{\genfunc}(a)
        \]
        Clearly, the function $\Gamma(\genpoly_1(a),...,\genpoly_c(a))$ is $\genpolyset$-measurable in $\field$,
        and because we have $\factor \semrefine \factor_1$, it is also $\genpolyset_1$-measurable in $\field$.
        Thus there exists $\funcdef{\Gamma_1}{\basefield^{c_1}}{\basefield}$ such that:
        \[
            \forall a \in \field: \genfunc(a) = \Gamma_1(\genpoly^{1}_{1}(a),...,\genpoly^{1}_{c_1}(a)) + \relativeremainder{\genfunc}(a)
        \]
        And therefore we have $\factor_1 \relsemrefineex{\variety} \factor$.
        \newline
        Now, we prove $\factor_2 \relsemrefineex{\variety} \factor_1$.
        Let $\funcdef{\genfunc}{\field}{\basefield}$ be a function that is $\genpolyset_1$-measurable relative to $\variety$.
        Again, we denote $\degree_{\genfunc} \definedas \deg(\genfunc)$,
        and by definition there exists $\funcdef{\Gamma_1}{\basefield^c}{\basefield}$,
        $\funcdef{\relativeremainder{\genfunc_1}}{\field}{\basefield}$ where $\deg(\genfunc - \relativeremainder{\genfunc_1}) \leq \degree_{\genfunc}$ and $\restrictfunc{\relativeremainder{\genfunc_1}}{\variety} \equiv 0$, such that:
        \begin{equation} \label{eq:f-decomposition-a}
            \forall a \in \field: \genfunc(a) = \Gamma_1(\genpoly^{1}_{1}(a),...,\genpoly^{1}_{c_1}(a)) + \relativeremainder{\genfunc}_1(a)
        \end{equation}
        Note that we also have $\deg(\relativeremainder{\genfunc_1}) \leq \degree_{\genfunc}$.
        We will refer this equation, and its simplifications we do throughout the proof, as \emph{the $\genpolyset_1$-decomposition of $\genfunc$}.
        \newline
        We wish to show that there exists $\funcdef{\Gamma_2}{\basefield^{c_2}}{\basefield}$ and
        $\funcdef{\relativeremainder{\genfunc}_2}{\field}{\basefield}$ where $\deg(\genfunc - \relativeremainder{\genfunc}_2) \leq \degree_{\genfunc}$ and $\restrictfunc{\relativeremainder{\genfunc}_2}{\variety} \equiv 0$, such that:
        \[
            \forall a \in \field: \genfunc(a) = \Gamma_2 \parens {\genpoly^1_1(a),...\genpoly^1_{i^\star - 1}(a), \genpoly^1_{i^\star + 1}(a),..., \genpoly^1_c(a), \genpoly[2]_1(a),...,\genpoly[2]_{\rankfunc(c_1)}(a)} + \relativeremainder{\genfunc}_2(a)
        \]
        We will do so using the $\genpolyset_1$-decomposition of $\genfunc$.
        Note that showing $\deg(\genfunc - \relativeremainder{\genfunc}) \leq \degree_{\genfunc}$ is equivalent of showing $\deg(\relativeremainder{\genfunc}_2) \leq \degree_{\genfunc}$.
        \newline
        First, by the way we built $\genpolyset_2$, using the same notations in the regularization step, we have:
        \[
            \forall a \in \field: \genpoly_{i^\star}^1(a) =
                \vec{\Gamma} \parens {\genpoly[2]_1(a),...,\genpoly[2]_{\rankfunc(c_1)}(a)}
                + \relativeremainder{\genpoly}(a)
                - \sum_{i \neq i^\star}{\genpoly_{i}^1(a)}
        \]
        Next, we substitute the value of $\genpoly_{i^\star}^1$ in the $\genpolyset_1$-decomposition of $\genfunc$ (\ref{eq:f-decomposition-a}),
        and get another decomposition of $\genfunc$ that does not depend on $\genpoly_{i^\star}^1$.
        Specifically we have:
        \begin{align} \label{eq:f-decomposition-b}
            \forall a \in \field: \genfunc(a) &=
        \Gamma_1 \parens
            {\genpoly^{1}_{1}(a)
                ,...,
                \parens{
                    \vec{\Gamma} \parens {\genpoly[2]_1(a),...,\genpoly[2]_{\rankfunc(c_1)}(a)}
                    + \relativeremainder{\genpoly}(a)
                    - \sum_{i \neq i^\star}{\genpoly_{i}^1(a)}}
                ,...,
                \genpoly^{1}_{c_1}(a))}\\
            &+ \relativeremainder{\genfunc}_1(a)
        \end{align}
        We wish to use the equation above to show that $\genfunc$ is $\genpolyset_2$-measurable relative to $\variety$.
        However, in order to show that the equation above is in the desired structure that proves that $\genfunc$ is $\genpolyset_2$-measurable,
        the expression inside $\Gamma_1$ must not depend on $\relativeremainder{\genpoly}$ because $\relativeremainder{\genpoly} \notin \genpolyset_2$.
        Note that this is enough as the rest of the polynomials in the expression above are in $\genpolyset_2$,
        and therefore without $\relativeremainder{\genpoly}$ the expression is $\genpolyset_2$-measurable.
        \newline
        To do so, we start by simplifying some of the notations.
        We denote:
        \[
            \vec{\genpoly}_2(a) \definedas \vec{\Gamma} \parens {\genpoly[2]_1(a),...,\genpoly[2]_{\rankfunc(c_1)}(a)} - \sum_{i \neq i^\star}{\genpoly_{i}^1(a)}
        \]
        This is the part of the sum that decomposes $\genpoly_{i^\star}^1(a)$ that is $\genpolyset_2$-measurable,
        thus the following equality applies:
        \[
            \genpoly_{i^\star}^1(a) = \vec{\genpoly}_2(a) + \relativeremainder{\genpoly}(a)
        \]
        where $\deg(\vec{\genpoly_2}), \deg(\relativeremainder{\genpoly}) \leq \degree^{\star}$.
        Using this notation, we write the $\genpolyset_1$-decomposition of $\genfunc$ (\ref{eq:f-decomposition-b}), and get:
        \begin{equation} \label{eq:f-decomposition-c}
            \forall a \in \field: \genfunc(a) = \Gamma_1 \parens
            {\genpoly^{1}_{1}(a)
                ,...,
                \parens{
                    \vec{\genpoly}_2(a)
                    - \relativeremainder{\genpoly}(a)}
                ,...,
                \genpoly^{1}_{c_1}(a))}
            + \relativeremainder{\genfunc}_1(a)
        \end{equation}
        Now, we use the following key observation:
        $\rank{\genpolyset_1} \geq \rankfunc_{\ref{preserving-degree-starting-field}}(\basefield, \degree, c_1)$,
        and as $\deg(\Gamma_1(\genpoly_1^1,...,\genpoly^1_{c_1})) \leq \degree_{\genfunc}$
        we can use Lemma~\ref{preserving-degree-starting-field} to achieve that $\Gamma_1$ is a polynomial of the form:
        \begin{equation*} \label{eq:regularization-gamma-is-a-polynomial}
            \Gamma_1(z_1,...,z_{c_1})  =
            \sum_{\alpha \in \sparens{\basefieldsize - 1}^{c_1}}
            {C_{\alpha} \cdot {\prod_{i = 1}^{c_1}}{z_i^{\alpha_i}}}
        \tag{$\star$}
        \end{equation*}
        where $C_{\alpha} = 0$ whenever $\sum_{i = 1}^{c_1}(\alpha_i \cdot \deg(\genpoly_i^1)) > \degree_{\genfunc}$.
        \newline
        Next, we substitute the polynomial structure of $\Gamma_1$ \eqref{eq:regularization-gamma-is-a-polynomial}
        in the $\genpolyset_1$-decomposition of $\genfunc$~\eqref{eq:f-decomposition-c},
        and observe what happens to each summand monomial with non-zero coefficients of $\Gamma_1$ in the expression after the substitution.
        \newline
        We will show that each such monomial is either $\genpolyset_2$-measurable,
        or a sum of a $\genpolyset_2$-measurable function with a valid $\variety$-remainder, i.e. a polynomial of degree $\leq \degree_{\genfunc}$ that is $\equiv 0$ in $\variety$.
        Note that if this is true for each monomial,
        every linear combination of such monomials is also a sum of $\genpolyset_2$-measurable function with a valid $\variety$-remainder.
        Thus, this will also be true for the entire decomposition of $\genfunc$, as it is a linear combination of such monomials summed with a valid remainder $\relativeremainder{\genfunc}_1$.
        This will conclude the proof.
        \newline
        Let $\alpha = (\alpha_1,...,\alpha_{c_1})$ be a vector of degrees that represents such a monomial.
        If $\alpha_{i^\star} = 0$, then the monomial is in the form:
        \[
            \prod_{i \in [c_1]}{{\genpoly_i}^{\alpha_i}} =
            \prod_{i \in [c_1] \setminus \set{i^{\star}}}{{\genpoly_i}^{\alpha_i}}
        \]
        and therefore it is clearly $\genpolyset_2$-measurable as all the polynomials in the expression above are in $\genpolyset_2$.
        \newline
        Next, if $\alpha_{i^\star} \neq 0$, then the monomial is in the form:
       \begin{equation} \label{eq:monomial-of-gamma}
            \prod_{i \in [c_1]}{{\genpoly_i}^{\alpha_i}} =
            (\vec{\genpoly_2} + \relativeremainder{\genpoly})^{\alpha_{i^\star}} \cdot
                \parens{\prod_{i \in [c_1] \setminus \set{i^{\star}}}{{\genpoly_i}^{\alpha_i}}}
       \end{equation}
        where $\sum_{i \in [c_1]}(\alpha_i \cdot \deg(\genpoly_i^1)) \leq \degree_{\genfunc}$.
        As $\deg(\vec{\genpoly_2} + \relativeremainder{\genpoly}) = \deg(\genpoly_{i^\star}) =\degree^{\star}$, we have:
        \[
            \deg \parens{\prod_{i \in [c_1] \setminus \set{i^{\star}}}{{\genpoly_i}^{\alpha_i}}} =
                \sum_{i \in [c_1] \setminus {i^{\star}}}(\alpha_i \cdot \deg(\genpoly_i^1))
                \leq \degree_{\genfunc} - \alpha_{i^\star} \cdot \degree^{\star}
        \]
        Now, we open the left brackets in (\ref{eq:monomial-of-gamma}), i.e $(\vec{\genpoly_2} + \relativeremainder{\genpoly})^{\alpha_{i^\star}}$.
        This enables us to separate the monomial to the part that only depend on $\vec{\genpoly_2}$ summed with a polynomial with bounded degree multiplied by $\relativeremainder{\genpoly}$ (and therefore a valid remainder).
        To be more specific, the monomial is in the form:
        \[
            (\vec{\genpoly_2} + \relativeremainder{\genpoly})^{\alpha_{i^\star}} = \vec{\genpoly_2}^{\alpha_{i^\star}} + \relativeremainder{\genpoly_{\alpha}}
        \]
        for some polynomial $\relativeremainder{\genpoly_{\alpha}}$ such that:
        \begin{enumerate}
            \item $\relativeremainder{\genpoly_{\alpha}}$ is of degree
                    $\deg(\relativeremainder{\genpoly_{\alpha}}) \leq \max \set{\deg(\vec{\genpoly_2}), \deg(\relativeremainder{\genpoly)}} \cdot \alpha_{i^\star} \leq \alpha_{i^\star} \cdot \degree^{\star}$
            \item $\relativeremainder{\genpoly_{\alpha}}$ is a multiple of $\relativeremainder{\genpoly}$, and therefore $\restrictfunc{\relativeremainder{\genpoly_{\alpha}}}{\variety} \equiv 0$
        \end{enumerate}
        Therefore, by substituting the left brackets back to the equation (\ref{eq:monomial-of-gamma}) and as $\vec{\genpoly_2}$ and $\genpoly_i$ for $i \neq i^{\star}$ are $\genpolyset_2$-measurable,
        one can see that the monomial is a sum of a $\genpolyset_2$-measurable polynomial with a valid remainder.
        Specifically, the remainder $\equiv 0$ in $\variety$, and its degree is $\leq \alpha_{i^\star} \cdot \degree^{\star} + \degree_{\genfunc} - \alpha_{i^\star} \cdot \degree^{\star} = \degree_{\genfunc}$.
        This concludes the proof of the claim.
        \end{proof}
    Now, it remains to prove the second part of the Theorem~{\ref{theorem:regularization-in-X}}.
    \begin{claim}
        If $\factor \synrefine \bar{\factor}$ for some polynomial factor $\bar{\factor}$ with
        relative rank of at least $\rankfunc(c^\prime)+c^\prime+1$ and rank of at least ${\rankfunc_{\ref{preserving-degree-starting-field}}(\basefield, \degree, c^{\prime})} + c^\prime + 1$,
        then we can require that $\factor^\prime \synrefine \bar{\factor}$.
    \end{claim}
    \begin{proof}
        We will show claim step-by-step.
        We denote by $\genpolyset, \bar{\genpolyset}, \genpolyset_1, \genpolyset_2$ the polynomial sets that generate the factors $\factor, \bar{\factor}, \factor_1, \factor_2$.
        Note that $\factor_1, \factor_2$ are the factors in the current step of the regularization process, and thus change in each step of the proof.
        We show that in each step, if $\factor \synrefine \bar{\factor}$ for some polynomial factor $\bar{\factor}$ with
        relative rank of at least $\rankfunc(c^\prime)+c^\prime+1$ and rank of at least ${\rankfunc_{\ref{preserving-degree-starting-field}}(\basefield, \degree, c^{\prime})} + c^\prime + 1$,
        then we can require that $\factor_1 \synrefine \bar{\factor}$, and also that $\factor_2 \synrefine \bar{\factor}$.
        \newline
        For the first part, we have $\factor_1 \synrefine \bar{\factor}$ by a simple usage of the second part of lemma~\ref{regularization-in-Fn-lemma},
        as:
        \[
            \rank{\bar{\genpolyset}}
            > \rankfunc_{\ref{preserving-degree-starting-field}}(\basefield, \degree, c^{\prime}) + c^\prime + 1
            \geq \rankfunc_{\ref{preserving-degree-starting-field}}(\basefield, \degree, c_1) + c_1 + 1
        \]
        \newline
        Now we prove the second part.
        We show that in the current regularization step, we could replace $\genpoly^1_{i^\star} \in \genpolyset_1$ such that $\genpoly^1_{i^\star} \notin \bar{\genpolyset}$.
        Note that this is possible whenever $\genpolyset^\star\cap \bar{\genpolyset} \neq \emptyset$ as the choice of $i^\star$ is arbitrary in polynomials which are in $\genpolyset^\star$.
        \newline
        Assume that is not possible and the factor $\genpolyset_1$ is still not $\rankfunc$-$\variety$-regular.
        Then, we have a linear combination $\vec{\genpoly}(x) \definedas \sum_{i=0}^{c_1}{\lambda_{i}\genpoly^1_i(x)}$ with $\drelrank{\degree^\star}{\variety}{\vec{\genpoly}} \leq \rankfunc(c_1)$
        where $\degree^\star = \max_{i \in \sparens{c_1}} {\deg(\lambda_i \genpoly^1_i)}$.
        We denote by $I^{\star} \subseteq [c_1]$ the set of indexes of such maximal-degree polynomials.
        By this notation, our assumption states that for all $i \in I^{\star}$ we have $\genpoly^1_i \in \bar{\genpolyset}$.
        Additionally, note that for all $i \notin I^{\star}$ we have $\deg(\genpoly^1_i) < \degree^{\star}$.
        Therefore, as the linear combination is of $\degree^\star$-relative rank $\leq \rankfunc(c_1)$,
        there exists a polynomial $\relativeremainder{\genpoly}$ of degree $\leq \deg(\vec{\genpoly}) \leq \degree^\star$ with $\restrictfunc{\relativeremainder{\genpoly}}{\variety} \equiv 0$
        such that $\drank{\degree^\star}{\vec{\genpoly} - \relativeremainder{\genpoly}} \leq \rankfunc(c_1)$.
        In other words, there exist a measurement function $\funcdef{\vec{\Gamma}}{\basefield^{\rankfunc(c_1)}}{\basefield}$ and polynomials $\genpoly[2]_1,...,\genpoly[2]_{\rankfunc(c_1)}$
        with $\deg(\genpoly[2]_i) \leq \degree^\star$ such that:
        \[
            \forall a \in \field: \vec{\genpoly}(a) - \relativeremainder{\genpoly}(a) = \vec{\Gamma} \parens {\genpoly[2]_1(a),...,\genpoly[2]_{\rankfunc(c_1)}(a)}
        \]
        By a simple calculation we have:
        \[
            \forall a \in \field: \sum_{i \in I^{\star}}{\genpoly^1_i(a)} - \relativeremainder{\genpoly}(a) =  \vec{\Gamma} \parens {\genpoly[2]_1(a),...,\genpoly[2]_{\rankfunc(c_1)}(a)} +  \sum_{i \notin I^\star}{\genpoly^1_i(a)}
        \]
        and by this we found a linear combination of polynomials in $\bar{\genpolyset}$ with maximal degree $\degree^\star$,
        that has $\degree^\star$-relative-rank $\leq \rankfunc(c_1) + c_1 + 1$.
        This is a contradiction to our assumptions on $\bar{\factor}$, which completes the proof of the claim.
    \end{proof}
    This completes the proof of the lemma.
\end{proof}

\section[Radius of Reed-Muller over \titlevariety]{Radius of Reed-Muller over \titlevariety}\label{sec:radius-of-RM-over-X}
We recall that the normalized distances of Reed-Muller codes over $\field$ and over $\variety$
are denoted by $\normalizedcodedistance{\basefield}{\degree}$ and $\normalizedcodedistanceex{\basefield}{\variety}{\degree}$ respectively.
We present a theorem that shows that Reed-Muller codes over a subset $\variety \subseteq \field$ that is $\degree$-lift-enabler and has the (limited) relative rank-bias property,
has (approximately) an \emph{equal} normalized distance as Reed-Muller codes over $\field$.
%TODO: Change to +- epsilon instead of >

\begin{theorem}\label{thm:distance-of-RM-in-X}
    There exist a function $\epsilon_1(\basefield, \degree, \rankbiasfunc, \epsilonlimitedrankbias)$  such that the following holds:
    Let $\basefield$ be a finite field, and let $\degree \in \naturalnumbersset$ be an integer that represents a degree.
    Let $\epsilonlimitedrankbias > 0$, and let $\funcdef{\rankbiasfunc}{[\epsilonlimitedrankbias, \infty]}{\naturalnumbersset}$ be a limited-relative-rank-bias function.
    \newline
    Let $\variety \subseteq \field$ be a set with the following properties
    \begin{enumerate}
        \item $\variety$ is $\degree$-lift-enabler with a lift operator $\lift{\square}$.
        \item $\variety$ has the $(\rankbiasfunc, \basefield, \degree, \epsilonlimitedrankbias)$-relative-rank-bias property.
    \end{enumerate}
    Then, for $\epsilon_1 \definedas \epsilon_1(\basefield, \degree, \rankbiasfunc, \epsilonlimitedrankbias)$ we have that for all $\blocklength \in \naturalnumbersset$:
    \[
        \normalizedcodedistanceex{\basefield}{\variety}{\degree} \geq \
        \normalizedcodedistance{\basefield}{\degree} - \epsilon_1
    \]
\end{theorem}

\begin{proof}
    We wish to do a reduction of our question regarding the radius of Reed-Muller in $\variety$ to the same question about Reed-Muller in $\field$.
    Let $\basefield$ be a finite field, and let $\degree \in \naturalnumbersset$ be an integer that represents a degree.
    Let $\epsilonlimitedrankbias > 0$, and let $\funcdef{\rankbiasfunc}{[\epsilonlimitedrankbias, \infty]}{\naturalnumbersset}$ be a limited-rank-relative-bias function.
    Let $\epsilon_1 \definedas \epsilon_1(\basefield, \degree, \rankbiasfunc, \epsilonlimitedrankbias)$ be a function we will specify later.
    Let $\variety \subseteq \field$ be a set with the properties defined above.
    \newline
    Moreover, let $\epsilon > \epsilon_1$ be some positive value.
    We will show that:
    \[
        \normalizedcodedistanceex{\basefield}{\variety}{\degree} > \
        \normalizedcodedistance{\basefield}{\degree} - \epsilon
    \]
    This will be enough as if the above holds for every $\epsilon > \epsilon_1$, we get that in fact
    $\normalizedcodedistanceex{\basefield}{\variety}{\degree} \geq \normalizedcodedistance{\basefield}{\degree} - \epsilon_1$.
    \newline
    For start, we note a simple observation: as Reed-Muller over $\variety$ is a linear code, we have
    \[
        \normalizedcodedistanceex{\basefield}{\variety}{\degree} =
        \min \set{\prex{x \in \variety}{\onvarpoly(x) \neq 0} \suchthat \onvarpoly \in \allpolyset{\leq \degree}{\variety}{\basefield}}
    \]
    Now, let $\onvarpoly \in \allpolyset{\leq \degree}{\variety}{\basefield}$ be a polynomial over $\variety$,
    and denote $\degree_{\onvarpoly} \definedas \deg(\onvarpoly)$.
    We wish to lower-bound the value of $\prex{x \in \variety}{\onvarpoly(x) \neq 0}$.
    To do so, we will equivalently upper-bound the value of $\prex{x \in \variety}{\onvarpoly(x) = 0}$.
    Precisely, to complete the proof all we need to show is:
    \[
        \prex{x \in \variety}{\onvarpoly(x) = 0} \leq 1 - \normalizedcodedistance{\basefield}{\degree} + \epsilon
    \]
    \newline
    Now we begin the proof itself.
    First, we lift the polynomial $\onvarpoly$ and get a polynomial $\funcdef{\lift{\onvarpoly}}{\field}{\basefield}$
    such that $\restrictfunc{\lift{\onvarpoly}}{\variety} \equiv \onvarpoly$ and $\deg(\lift{\onvarpoly}) = \degree_{\onvarpoly}$.
    Next, denote by $\factor_{\lift{\onvarpoly}}$ the factor defined by the set of single polynomial $\genpolyset = \set {\lift{\onvarpoly}}$.
    Trivially, the polynomial $\lift{\onvarpoly}$ is measurable in respect of $\genpolyset$.
    \newline
    We define the rank function:
    \[
        \rankfunc(m) \definedas \max \set{
            \rankbiasfunc \parens {\dfrac{\epsilon / 2}{\abs{\basefield}^m}},
            \rankfunc_{\ref{high-rank-implies-low-bias}} \parens{\basefield, \degree, \dfrac{\epsilon / 2}{\abs{\basefield}^m}}}
    \]
    Then, we $\rankfunc$-$\variety$-regularize $\genpolyset$ using Lemma~{\ref{theorem:regularization-in-X}}.
    This gives us a $\rankfunc$-$\variety$-regular factor $\factor^\prime$, which is defined by a set of polynomials $\genpolyset^{\prime} \definedas \set{\genpoly^\prime_1,...,\genpoly^\prime_{c^\prime}}$
    of degree $\leq \degree$ such that $\factor^\prime \relsemrefine{\variety}\factor_{\lift{\onvarpoly}}$ with $\relrank{\variety}{\genpolyset^\prime} \geq \rankfunc$
    and with bounded amount of polynomials defining it i.e,$c^\prime \leq C_{\rankfunc, \degree}^{\ref{theorem:regularization-in-X}}(1)$.
    Therefore, from definition we have that $\lift{\onvarpoly}$ is $\genpolyset^\prime$-measurable relative to $\variety$.
    Thus, there exists a measurement function $\funcdef{\Gamma}{\basefield^{c^\prime}}{\basefield}$
    and a remainder $\funcdef{\relativeremainder{\Gamma}}{\field}{\basefield}$ with $\restrictfunc{\relativeremainder{\Gamma}}{\variety} \equiv 0$
    and degree bounded by $\degree_{\onvarpoly}$, such that:
    \[
        \forall a \in \field:
        \lift{\onvarpoly}(a) =
        \Gamma(\genpoly^\prime_1(a),...,\genpoly^\prime_{c^\prime}(a))
        + \relativeremainder{\Gamma}(a)
    \]
    Next, we denote $\genpoly^\prime \definedas \lift{\onvarpoly} - \relativeremainder{\Gamma}$.
    By definition of remainder function, we have that $\restrictfunc{\genpoly^\prime}{\variety} \equiv \onvarpoly$.
    Additionally, note that $\genpoly^\prime$ is a polynomial over $\field$ of degree $\deg(\genpoly^\prime) = \degree_{\onvarpoly} \leq \degree$,
    and hence by the definition of $\normalizedcodedistance{\basefield}{\degree}$:
    \begin{equation}\label{eq:polynomials-in-fn-are-bounded-away-from-zero-with-high-probability}
    \prex{a \in \field}{\genpoly^\prime(a) = 0} \leq 1 - \normalizedcodedistance{\basefield}{\degree}
    \end{equation}
    For the next step, we claim that $\genpoly^\prime$ equals $0$ in $\field$ approximately with the same probability it equals $0$ in $\variety$.
    Note that this is the heart of the proof: it allows use properties known in $\field$ to new properties in $\variety$.
    This is formulated as follows:
    \begin{claim}\label{claim:p-and-P-have-the-same-approximation}
        We have:
        \[
            \abs{\prex{a \in \field}{\genpoly^\prime(a)= 0} -
            \prex{x \in \variety}{\genpoly^\prime(x) = 0}} \leq  \epsilon
        \]
    \end{claim}
    \begin{proof}
        Denote $S \definedas \basefield^{c^\prime}$, and for all $s \in S$, denote:
        \[
            p_1(s) \definedas \prex{a \in \field}{(\genpoly^{\prime}_1(a),...,\genpoly^{\prime}_{c^\prime}(a)) = s}
        \]
        As of our choice of $\rankfunc$, we have $\rank{\genpolyset^{\prime}} \geq \rankfunc_{\ref{high-rank-implies-low-bias}} \parens {\basefield, \degree, \dfrac{\epsilon/2}{\abs{\basefield}^{c^\prime}}}$.
        By combining Lemma~{\ref{high-rank-implies-low-bias}} with Lemma~{\ref{every-linear-combination-has-low-bias-implies-equidistribution}},
        we have that $p_1$ is ($\epsilon/2\abs{S}$)-equidistributed, i.e:
        \[
            p_1(s) = \dfrac{1 \pm \epsilon/2}{\abs{S}}
        \]
        Similarly, denote:
        \[
            p_2(s) \definedas \prex{x \in \variety}{(\genpoly^{\prime}_1(x),...,\genpoly^{\prime}_{c^\prime}(x)) = s}
        \]
        As of our choice of $\rankfunc$, we have $\relrank{\variety}{\genpolyset^\prime} \geq \rankbiasfunc(\epsilon / 2\abs{S})$.
        Now, we wish to use the relative rank-bias relation with Lemma~\ref{every-linear-combination-has-low-bias-implies-equidistribution}
        to conclude similarly that $p_2$ is ($\epsilon/2\abs{S}$)-equidistributed, i.e:
        \[
            p_2(s) = \dfrac{1 \pm \epsilon/2}{\abs{S}}
        \]
        However, in order to so, we must first ensure that $(\epsilon/2\abs{S}) \geq \epsilonlimitedrankbias$.
        This is done by choosing a correct $\epsilon_1$, and formulated in the following claim:
        \begin{claim}
            One can choose $\epsilon_1 \definedas \epsilon_1(\basefield, \degree, \rankbiasfunc, \epsilonlimitedrankbias)$ such that if $\epsilon \geq \epsilon_1$ we have that $\epsilon/2\abs{S} \geq \epsilon_1$.
        \end{claim}
        \begin{proof}
            We need that:
            \[
                \dfrac{\epsilon}{2 \abs{\basefield}^{c^\prime}} \geq \epsilonlimitedrankbias
            \]
            As $c^\prime \leq C^{\ref{theorem:regularization-in-X}}_{\rankfunc, \degree}(1)$,
            for the term above to hold it is enough that the following will be true:
            \[
                \epsilon \geq \epsilonlimitedrankbias \cdot 2 \abs{\basefield}^{C^{\ref{theorem:regularization-in-X}}_{\rankfunc, \degree}(1)}
            \]
            and as $\rankfunc$ and thus also $C^{\ref{theorem:regularization-in-X}}_{\rankfunc, \degree}(1)$ are independent of $\blocklength$,
            we can pick $\epsilon_1 = \epsilon_1(\basefield, \degree, \rankbiasfunc, \epsilonlimitedrankbias)$ and get what we aimed for.
        \end{proof}
        Now, under that assumption of $\epsilon_1$ written above, we have that $p_2$ is ($\epsilon/2\abs{S}$)-equidistributed.
        This allows us to use the similar distributions of $\genpolyset^\prime$ in $\field$ and in $\variety$ to conclude
        that $\genpoly^\prime$ behaves similar in $\field$ and in $\variety$:
        \begin{flalign*}
            \prex{a \in \field}{\genpoly^\prime(a)= 0}
            &=\sum_{s \in S} {p_1(s) \cdot \existfunc{\Gamma(s) = 0}} \\
            &=\sum_{s \in S} {p_2(s) \cdot \existfunc{\Gamma(s) = 0}} \pm \epsilon \\
            &=\prex{x \in \variety}{\genpoly^\prime(x)= 0} \pm \epsilon
        \end{flalign*}
        which concludes the proof of the claim.
    \end{proof}

    Finally, as $\restrictfunc{\genpoly^\prime}{\variety} \equiv \onvarpoly$, we have that $\prex{x \in \variety}{\genpoly^\prime(x) = 0} = \prex{x \in \variety}{\onvarpoly(x) = 0}$.
    Thus, the claim above combining with~\eqref{eq:polynomials-in-fn-are-bounded-away-from-zero-with-high-probability}
    shows that the probability we wished to bound is bounded as we aimed for:
    \[
        \prex{x \in \variety}{\onvarpoly(x) = 0} \leq 1 - \normalizedcodedistance{\basefield}{\degree} + \epsilon
    \]
    This concludes the proof of the theorem.
\end{proof}
\begin{remark}
    Under the same conditions,
    the distance of Reed-Muller codes in $\variety$ is also bounded \emph{from above} by the distance of Reed-Muller codes in $\field$,
    and we have:
    \[
        \normalizedcodedistanceex{\basefield}{\variety}{\degree} \leq \normalizedcodedistance{\basefield}{\degree} + \epsilon_1
    \]
    \begin{proof}
        Let $\funcdef{\genfunc}{\field}{\basefield}$ be the polynomial in $\field$ with the \emph{smallest} distance from $0$ as possible, that is $\normalizedcodedistance{\basefield}{\degree}$.
        Denote $\onvarpoly \definedas \restrictfunc{\genpoly}{\variety}$.
        Note that $\onvarpoly$ is a polynomial in $\variety$.
        Now repeat the proof using these two polynomials, and by Claim~\ref{claim:p-and-P-have-the-same-approximation}, we have that a random input of $\genpoly$ yields $0$
        (approximately) the same as a random input of $\onvarpoly$ yields $0$.
        Thus as we have $\prex{x \in \field}{\genpoly(x) = 0} = 1 - \normalizedcodedistance{\basefield}{\degree}$
        we also get:
        \[
            \prex{x \in \variety}{\onvarpoly(x) = 0} \geq 1 - \normalizedcodedistance{\basefield}{\degree} - \epsilon_1
        \]
        This bounds \emph{from above} the distance of Reed-Muller code in $\variety$ and we have:
        \[
            \normalizedcodedistanceex{\basefield}{\variety}{\degree} \leq \normalizedcodedistance{\basefield}{\degree} + \epsilon_1
        \]
    \end{proof}
\end{remark}
\begin{corollary}
    If we assume $\variety$ has the limited-relative rank-bias property to \emph{any extent} (or just the relative rank-bias property),
    then the theorem above proves an exact equality $\normalizedcodedistanceex{\basefield}{\variety}{\degree} = \normalizedcodedistance{\basefield}{\degree}$.
\end{corollary}

\section[Equidistribution of Functions]{Equidistribution of Functions}\label{sec:equidistribution-of-functions}
Assume we have a collection of functions $(\genfunc_1,...,\genfunc_c$), where $\funcdef{\genfunc_i}{A}{\basefield}$ for some finite set $A$.
We are interested in showing that the functions are equidistributed, which means that their values behave close to independent random variables.
We begin by formulating this definition:
\begin{definition}[Equidistribution of Functions]
    Given $\epsilon > 0$ and $A \subseteq \field$,
    we say a collection of functions $\genfuncset = (\genfunc_1,...,\genfunc_c)$ where $\funcdef{\genfunc_i}{\genset}{\basefield}$ is $\epsilon$-equidistributed in $A$ if for all $\vec{\alpha} = (\alpha_1,...,\alpha_c) \in \basefield^c$ we have:
    \[
        \prex{x \in A}{(\genfunc_1(x),...,\genfunc_c(x)) = \vec{\alpha}} = \frac{1}{\abs{\basefield}^c} \pm \epsilon
    \]
\end{definition}

The following is a standard lemma that shows that if every linear combination of a collection of functions has low bias, the collection is equidistributed.
We repeat the steps of the proof of \cite[Lemma 7.24]{book}, but here, we think of $A$ as any finite set (and not particularly $\field$):
\begin{lemma}\label{every-linear-combination-has-low-bias-implies-equidistribution}
Let $\epsilon > 0$, and let $A$ be a finite set.
Let $\genfuncset = (\genfunc_1,...,\genfunc_c)$ be a collection of functions defined over $A$, i.e. $\funcdef{\genfunc_i}{A}{\basefield}$.
Assume each linear combination of the collection has low bias, i.e for each $\lambda = (\lambda_1,...,\lambda_c) \in \basefield^c$ such that $\lambda \neq \vec{0}$ we have:
\[
    \relbias{x \in A}{\sum_{i=1}^{c}{\lambda_i \genfunc_i}} < \epsilon
\]
Then, the collection $\genfuncset$ is $\epsilon$-equidistributed over $A$.
\newline
In particular, for $\epsilon < \frac{1}{\abs{\basefield}^c}$, the lemma shows that each atom of $\genfuncset$ is not empty i.e for all $\vec{\alpha}$ there is some $x \in A$ such that $(\genfunc_1(x),...,\genfunc_c(x)) = \vec{\alpha}$.
\end{lemma}
\begin{proof}
    We wish to show that for each $\vec{\alpha} \in \basefield^c$ we have:
    \[
        \prex{x \in A}{(\genfunc_1(x),...,\genfunc_c(x)) = \vec{\alpha}} = \frac{1}{\abs{\basefield}} \pm \epsilon
    \]
    We express the fraction of inputs that are in the atom $\vec{\alpha}$ the following way:
    \[
        \prex{x \in A}{(\genfunc_1(x),...,\genfunc_c(x))} =
        \expectation{x \in A}{\prod_{i=1}^c {1_{[\genfunc_i(x) = \alpha_i]}}}
    \]
    We use the fact that for every $0 \neq x \in \basefield$, we have $\sum_{\lambda = 0}^{\basefieldsize - 1} \charfunc{\lambda x} = 0$,
    and if $x = 0$ we have $\sum_{\lambda = 0}^{\basefieldsize - 1} \charfunc{\lambda x} = \basefieldsize$.
    Therefore, the expression above equals:
    \[
        =\expectation{x \in A}{\prod_{i=1}^c \parens {{\frac{1}{\basefieldsize} \cdot \sum_{\lambda_i = 0}^{\basefieldsize - 1} {\charfunc{\lambda_i (\genfunc_i(x) - \alpha_i)}}}}} =\\
        \frac{1}{\basefieldsize^c} \cdot {\expectation{x \in A}{\prod_{i=1}^c \sum_{\lambda_i = 0}^{\basefieldsize - 1} {\charfunc{\lambda_i (\genfunc_i(x) - \alpha_i)}}}}
    \]
    By the definition of character functions, we have that $\charfunc{a+b} = \charfunc{a} \cdot \charfunc{b}$, and therefore the expression above equals:
    \[
        \frac{1}{\basefieldsize^c} \cdot \sum_{(\lambda_1,...,\lambda_c) \in \prod_{i=1}^c [0, \basefieldsize - 1]} \parens{\expectation{x \in A}{\charfunc{\sum_{i = 0}^{c} {\lambda_i (\genfunc_i(x) - \alpha_i)}}}}
    \]
    Now, we use the fact that:
    \[
        \relbias{x \in A}{\sum_{i=1}^c (\lambda_i (\genfunc_i(x) -\alpha_i)} = \relbias{x \in A}{\sum_{i=1}^c (\lambda_i \genfunc_i(x))} < \epsilon
    \]
    and get that:
    \[
        \prex{x \in A}{(\genfunc_1(x),...,\genfunc_c(x)) = \vec{\alpha}} = \frac{1}{\basefieldsize^c} \cdot \parens{1 \pm \epsilon \prod_{i=1}^c{\basefieldsize}} = \frac{1}{\abs{\basefield}^c} \pm \epsilon
    \]
\end{proof}

%Now, if we have $A = \basefield^{\blocklength}$ for some $\blocklength \in \naturalnumbersset$, we have the following corollary:
%\begin{corollary}[High rank implies equidistribution in $\field$]
%    TODO %TODO
%\end{corollary}
%
%Moreover, if we have that $A = \variety$ when $\variety \subseteq \field$ has the relative rank-bias property, we have:
%\begin{corollary}[High relative rank implies equidistribution in $\variety$]
%    TODO %TODO
%\end{corollary}
%


\section[Comparing Ranks]{Comparing Ranks}\label{sec:comparing-ranks}
In this section, we compare the definition of rank we used in this paper to another definition of rank used implicitly throughout this paper.
This comparison is crucial, as there is no universally accepted definition of rank;
different theorems presented throughout this paper employ distinct definitions.
We demonstrate that our definition is sufficiently comprehensive, in that a polynomial (or a factor) classified as having high rank according to our criteria
also exhibits high rank according to the second implicitly-used definition.
While in many cases the comparison may appear straightforward, we include it for the sake of completeness.
\newline
Specifically, we compare our definition of rank with the definition established in ~\cite{lampert2021relative}.
The paper~\cite{lampert2021relative} extended the original definition of rank that was presented in \cite{10.1007/BF02392473},
to include also the concept of relative rank.
It is important to note that this definition is specifically defined to subsets $\variety \subseteq \field$ that can be expressed as sets in the form $\variety = \zerofunc{\varpolyset}$ for some set of polynomials $\varpolyset$,
and not to a general set $\variety \subseteq \field$.
\newline
First, we present a useful notation that is used in the definition presented in ~\cite{lampert2021relative}:
\begin{notation*}[Largest Degree Homogenous Part]
    For a polynomial $\genpoly$ of degree $\degree$, we denote by $\homopart{\genpoly}$ its degree-$\degree$ homogenous component.
    In other words, $\homopart{\genpoly}$ is the sum of all the monomials of $\genpoly$ of degree exactly $\degree$.
    For a set of polynomials $\genpolyset = \set{\genpoly_1,...,\genpoly_c}$, we define $\homopart{\genpolyset} \definedas \set{\homopart{\genpoly_i} \suchthat i = 1,...,c}$.
\end{notation*}
Next, we present the exact definition of rank for a polynomial:
\begin{definition}[Schmidt Rank of a Polynomial]
    The schmidt rank of a homogenous polynomial $\funcdef{\genpoly[1]}{\field}{\basefield}$, noted as $\schmrank{\genpoly[1]}$, is the minimal $r$ such that there exist $(\genpoly[2]_i, \genpoly[3]_i)_{i\in[r]}$
    with $\deg{\genpoly[2]_i}, \deg{\genpoly[3]_i} < \deg{\genpoly[1]}$ such that:
    \[
        \genpoly[1](x) = \sum_{i=1}^{r}(\genpoly[2](x) \cdot \genpoly[3](x))
    \]
    For a general polynomial $\genpoly$ of degree $\degree$, we set its rank to be the rank of its degree-$\degree$ homogenous component, i.e. $\schmrank{\genpoly} \definedas \schmrank{\homopart{\genpoly}}$.
\end{definition}

\begin{remark}[High rank implies high schmidt rank]\label{high-rank-implies-high-schmidt-rank}
If $\rank{\genpoly} \geq 2 \cdot r + 1$ for some constant $r \in \naturalnumbersset$, then $\schmrank{\genpoly} \geq r$.
\end{remark}
\begin{proof}
    For homogenous polynomial $\genpoly$, assume $\schmrank{\genpoly} < r$.
    Then, there exist $r^\prime < r$ such that there exist $(\genpoly[2]_i, \genpoly[3]_i)_{i=1}^{r^\prime}$ with $\deg{\genpoly[2]_i}, \deg{\genpoly[3]}_i < \deg \genpoly$ such that:
    \[
        \genpoly[1](x) = \sum_{i=1}^{r^\prime}(\genpoly[2](x) \cdot \genpoly[3](x))
    \]
    Then we can choose $\funcdef{\Gamma}{\basefield^{2r^\prime}}{\basefield}$ to be a sum of multiples of each two consecutive variables to get that
    $\genpoly(x) = \Gamma(\genpoly[2]_1(x), \genpoly[3]_1(x),...,\genpoly[2]_{r^\prime}(x), \genpoly[3]_{r^\prime}(x))$, where the polynomials are from a degree $< \deg(\genpoly)$.
    This means that $\rank{\genpoly} \leq 2r^{\prime} < 2r$ as we requested.
    \newline
    If we do not assume $\genpoly$ is homogenous, by adding $\genpoly - \homopart{\genpoly}$ as an input to $\Gamma$,
    one can create a $\funcdef{\Gamma^\prime}{\basefield^{2r^\prime+1}}{\basefield}$
    which equals to $\genpoly$ when substituting the inputs with some polynomials with degree $< \deg{\genpoly}$,
    which concludes the proof in a similar way.
\end{proof}

Next, we present the definition of Schmidt rank of a factor as defined in ~\cite{lampert2021relative}.
\begin{definition}[Schmidt Rank of a Factor]
    For a factor of homogenous polynomials $\genpolyset = (\genpoly_1,...,\genpoly_{c})$, the schmidt rank of the factor is defined as:
    \[
        \schmrank{\genpolyset} \definedas \min \parens{\schmrank{\sum_{i=1}^{c}\lambda_i \genpoly_i} \suchthat 0 \neq (\lambda_1,...,\lambda_c) \in \basefield^c}
    \]
    Similarly, for a factor of general polynomials $\genpolyset$, we set its rank to be the rank of its matching homogenous-factor,
    i.e. $\schmrank{\genpolyset} \definedas \schmrank{\homopart{\genpolyset}}$
    For a factor $\factor$ generated by $\genpolyset$, we define $\schmrank{\factor} \definedas \schmrank{\genpolyset}$.
\end{definition}

To establish the equivalence of this definition with the one employed throughout the paper, we must first acknowledge two key distinctions between the definitions.
The first distinction is that this definition focuses on the largest-degree homogeneous components of the polynomials involved in the factor, rather than considering linear combinations of polynomials from the factor.
The second distinction pertains to the treatment of $\degree$ in the computation of $\degree$-rank of each linear combination.
This definition uses the degree of the linear combination directly to calculate the rank that participates in the minimum, in contrast to our definition which uses $\max_{i}{\deg(\lambda_i \genpoly_i)}$.
Despite these differences, we will demonstrate that both definitions ultimately yield a similar rank assessment, thereby affirming their equivalence.
\begin{remark}[High Rank Implies High Schmidt Rank for Factors]
    Let $\genpolyset = \parens{\genpoly_1,...,\genpoly_c}$ be a set of polynomials and let $\rankval \in \naturalnumbersset$ be a positive integer, i.e. $\rankval > 0$.
    If $\rank{\genpolyset} \geq 2 \cdot \rankval + 1$, then $\schmrank{\genpolyset} \geq \rankval$.
    \begin{proof}
        Assume that $\schmrank{\genpolyset} \leq \rankval$ for $\rankval > 0$.
        We will show that $\rank{\genpolyset} \leq 2 \rankval + 1$.
        By definition, there exists a linear combination of polynomials in $\homopart{\genpolyset}$ with rank $\leq \rankval$.
        In other words, there exists $\vec{0} \neq \lambda \in \basefield^{c}$ such that $\schmrank{\sum_{i=1}^{c} {\lambda_i \homopart{\genpoly_i}}} \leq \rankval$.
        Denote $\vec{\genpoly_h} \definedas \sum_{i=1}^{c} {\lambda_i \homopart{\genpoly_i}} $.
        As was shown in a previous remark, a rank of a polynomial is smaller than its schmidt rank up to a constant factor,
        thus $\rank{\vec{\genpoly_h}} \leq 2\rankval + 1$ (see Remark~\ref{high-rank-implies-high-schmidt-rank}).
        \newline
        Next, we denote $\vec{\genpoly} \definedas \sum_{i=1}^{c} {\lambda_i \genpoly_i}$, and $\degree_M \definedas \max_{i \in \sparens{c}}{\lambda_i \genpoly_i}$.
        Note that  $\deg(\vec{\genpoly}) \leq \degree_M$.
        We wish to show that $\drank{\degree_M}{\vec{\genpoly}} \leq 2 \rankval + 1$.
        First, we observe that the $\degree_M$-degree homogenous component of $\vec{\genpoly_h}$ equals the $\degree_M$-degree homogenous component of $\vec{\genpoly}$.
        This is true because every highest-degree component of polynomials in the linear combination that generated $\vec{\genpoly}$,
        also exists in the linear combination that generates $\vec{\genpoly_h}$.
        In particular, all homogenous components of degree $\degree_M$ exists in both linear combinations $\vec{\genpoly_h}$ and $\vec{\genpoly}$.
        Therefore, if the degree of $\vec{\genpoly}$ equals $\degree_M$, we have that $\drank{\degree_M}{\vec{\genpoly}} = \rank{\vec{\genpoly}} \geq 2 \rankval + 1$.
        Otherwise, if $\deg(\vec{\genpoly}) < \degree_M$, then $\drank{\degree_M}{\vec{\genpoly}} = 1 \leq 2 \rankval + 1$.
        This completes the proof.
        \newline
    \end{proof}
\end{remark}
\begin{note*}
    In the case discussed above, if $\deg(\vec{\genpoly}) < \degree_M$, then $\schmrank{\genpolyset} = 0$.
    \begin{proof}
        Assume that $\deg(\vec{\genpoly}) < \degree_M$.
        Therefore, the degree of the linear combination $\vec{\genpoly} = \sum_{i=1}^{c} {\lambda_i \genpoly_i}$ is strictly smaller than the degree of at least one of the polynomials participating in it.
        Denote by $\vec{\lambda}^\star$ the sub-combination of $\vec{\lambda}$ that consists only the polynomials that participated in $\vec{\genpoly}$ that are of degree $= \degree_M$.
        Trivially, $\vec{\lambda}^\star \neq \vec{0}$.
        Additionally, we have $\deg(\sum_{i=1}^{c} {\lambda_i^\star \genpoly_i}) < \degree_M$.
        Now, we use the following observation: the linear combination above, when summing only the homogenous components of each polynomial, equals $0$, i.e. $\sum_{i=1}^{c} {\lambda_i^\star \homopart{\genpoly_i}} \equiv 0$.
        By this, we found a linear combination of $\homopart{\genpolyset}$ that is $\equiv 0$.
        Thus by definition, we have $\schmrank{\genpolyset} = 0$.
    \end{proof}
\end{note*}
\begin{note*}
    This shows that if we compare only the differences in the definition of rank of a factor, i.e. the focus on linear combinations of the largest-degree homogenous components in contrast to the use of the maximal degree $\degree$-rank,
    the two definitions for a rank of a factor are equal up to $\pm 1$
    (in case we use the same definition of rank for a single polynomial).
    To avoid confusion, we omit the exact definitions and respective proof.
\end{note*}

We now present the definition of relative rank as stated in ~\cite[Definition 1.6]{lampert2021relative}:
We remind the reader that this definition is specifically defined to subsets $\variety \subseteq \field$ that can be expressed by $\variety = \zerofunc{\varpolyset}$ for some set of polynomials $\varpolyset$, and not to a general set $\variety \subseteq \field$.
\begin{definition}[Relative Schmidt Rank of a Polynomial]
    The relative schmidt rank of a homogeneous polynomial $\genpoly[1]$ relative to a collection of homogeneous polynomials $\varpolyset=(\series{\varpoly}{1}{\varietypolycount})$ is
    \[
        \relschmrank{\varpolyset}{\genpoly[1]} \definedas
        \min \set{\schmrank{P+\sum_{i=1}^{\varietypolycount}{\remainderpoly_{i}\varpoly_{i}}}
            \suchthat
            \deg(\varpoly_{i})+\deg(\remainderpoly_{i}) \leq \deg(\genpoly[1]), \forall i \in [\varietypolycount]}
    \]
    Note that whenever $\deg{\varpoly_{i}}>\deg{\genpoly}$, this implies $\remainderpoly_{i}=0$.
    \newline
    For general polynomial $\genpoly$ and general collection of polynomials $\varpolyset$, we define
    the schmidt rank of the former in respect to the latter by the relative rank of their largest-degree homogenous component,
    i.e. $\relschmrank{\varpolyset}{\genpoly} \definedas \relschmrank{\homopart{\varpolyset}}{\homopart{\genpoly}}$.
\end{definition}
\begin{remark}[High Relative Rank $\Rightarrow$ High Relative Schmidt Rank]\label{remark-high-relative-rank-implies-high-relative-schmidt-rank}
    Let $\genpoly$ and $\varpolyset = \set{\varpoly_1,...,\varpoly_{\varietypolycount}}$ be polynomials,
    and let $\variety \subseteq \field$ be defined as $\variety = \zerofunc{\varpolyset}$.

    If $\relrank{\variety}{\genpoly} \geq 2 \cdot \rankval + 2$ for some constant $\rankval \in \naturalnumbersset$,
    then $\relschmrank{\varpolyset}{\genpoly} \geq r$.
\end{remark}
\begin{proof}
    Let $\genpoly$ and $\varpoly_1,...,\varpoly_{\varietypolycount}$ be polynomials.
    Assume that $\relschmrank{\varpolyset}{\genpoly} \leq \rankval$.
    Then, there exists $\remainderpoly_1,...,\remainderpoly_\varietypolycount$ with
    $\deg(\varpoly_{i})+\deg(\remainderpoly_{i}) \leq \deg(\genpoly[1])$ for all $i \in \sparens{\varietypolycount}$
    such that:
    \[
        \schmrank{\homopart{\genpoly} + \sum_{i=1}^{\varietypolycount}{\remainderpoly_i \homopart{\varpoly_i}}} \leq \rankval
    \]
    Denote $\relativeremainder{\genpoly_h} \definedas \sum_{i=1}^{\varietypolycount}{\remainderpoly_i \homopart{\varpoly_i}}$.
    As we have shown earlier, a rank of a polynomial is smaller than its schmidt rank up to a constant factor (See Remark~\ref{high-rank-implies-high-schmidt-rank}).
    Thus:
    \[
        \rank{\homopart{\genpoly} + \relativeremainder{\genpoly_h}} \leq
        2 \cdot \schmrank{\homopart{\genpoly} + \relativeremainder{\genpoly_h}} + 1 \leq
        2 \cdot \relschmrank{\variety}{\genpoly} + 1 =
        2 \rankval + 1
    \]
    Next, we denote the respective remainder polynomial for the non-homogenous analogue, i.e $\relativeremainder{\genpoly} \definedas \sum_{i=1}^{\varietypolycount}{\remainderpoly_i \varpoly_i}$.
    By observing the highest degree homogenous component of each summand, one can see that $\homopart{\genpoly + \relativeremainder{\genpoly}} = \homopart{\homopart{\genpoly} + \relativeremainder{\genpoly_h}}$.
    Therefore, by adding to the decomposition the non higest-degree-homogenous-component, one can see that:
    \[
        \rank{\genpoly + \relativeremainder{\genpoly}} \leq
        \rank{\homopart{\genpoly} + \relativeremainder{\genpoly_h}} + 1 \leq
        2 \rankval + 2
    \]
    This completes the proof as $\relrank{\variety}{\genpoly} \leq \rank{\genpoly + \relativeremainder{\genpoly}} \leq 2 \rankval + 2$.
\end{proof}

\begin{remark}[Relative Schmidt Rank over Varieties of High Degree]\label{relative-schimdt-rank-equals-schmidt-rank-if-the-variety-is-of-high-degree}
If the polynomials defining the variety $\varpolyset = (\varpoly_1,...,\varpoly_{\varietypolycount})$ are of degree $> \deg(\genpoly)$,
then, $\relschmrank{\varpolyset}{\genpoly} = \schmrank{\genpoly}$.
This is true because in this case, in the calculation of the minimum in the definition of relative schmidt rank, we must have $\remainderpoly_i = 1$ for all $i \in [\varietypolycount]$ and therefore the minimum above is simply $\rank{\genpoly}$.
\newline
Note that a similar statement holds for factors aswell.
If $\genpolyset = (\genpoly_1,...,\genpoly_c)$ is a factor of degree $\degree$, then if all the polynomials in $\varpolyset$ have degree $> \degree$, then the statement above is also true i.e $\relschmrank{\varpolyset}{\genpolyset} = \schmrank{\genpolyset}$.
This is true because for every linear combination of $\genpolyset$ has degree $\leq \degree$ and therefore its relative schimdt rank equals its rank.
\end{remark}

%TODO: Add a note, is the same property holds in our definition?

Finally, we present the extension of the definition of relative rank for polynomials factors:
\begin{definition}[Relative Schmidt Rank of a Factor]
    The relative rank of a set of homogenous polynomials $\genpolyset = \set{\genpoly_1,...,\genpoly_c}$
    relative to another collection of polynomials $\varpolyset = \set{\varpoly_1,...,\varpoly_\varietypolycount}$ is defined as:
    \[
        \relschmrank{\varpolyset}{\genpolyset} \definedas
        \min \set{\relschmrank{\varpolyset}{\sum_{i=1}^c{\lambda_i \genpoly_i}} \suchthat \vec{0} \neq (\lambda_1,...,\lambda_c) \in \basefield^c}
    \]
    If $\genpolyset$ is a general collection of polynomials, then $\relschmrank{\varpolyset}{\genpolyset} \definedas \relschmrank{\varpolyset}{\homopart{\genpolyset}}$.
    \newline
    For a factor $\factor$ generated by a set of polynomials $\genpolyset$, we define its schmidt rank relative to $\variety = \zerofunc{\varpolyset}$
    to be $\relschmrank{\variety}{\factor} \definedas \relschmrank{\varpolyset}{\genpolyset}$.
\end{definition}


\begin{remark}
    Let $\genpolyset = \set{\genpoly_1,...,\genpoly_c}$ and $\varpolyset = \set{\varpoly_1,...,\varpoly_{\varietypolycount}}$ be sets of polynomials,
    and let $\variety \subseteq \field$ be defined as $\variety = \zerofunc{\varpolyset}$.
    Additionally, let $\rankval \in \naturalnumbersset$ such that $\rankval > 0$.
    If $\relrank{\variety}{\genpolyset} \geq 2 \cdot \rankval + 2$ for some constant $\rankval \in \naturalnumbersset$, then $\relschmrank{\varpolyset}{\genpoly} \geq r$.
\end{remark}
\begin{proof}
    Assume that $\relschmrank{\varpolyset}{\genpolyset} \leq \rankval$.
    We will show that $\relrank{\variety}{\genpolyset} \leq 2\rankval + 2$.
    Let $\vec{0} \neq \vec{\lambda} \in \basefield^c$ be some vector of coefficients.
    Let $\vec{\genpoly} \definedas \sum_{i=1}^{c} {\lambda_i \genpoly_i}$ and $\vec{\genpoly_h} \definedas \sum_{i=1}^{c} {\lambda_i \homopart{\genpoly_i}}$
    be the linear combinations of polynomials in $\genpolyset$ and $\homopart{\genpolyset}$ with coefficients $\vec{\lambda}$ respectively,
    and let $\degree_M \definedas \max_{i \in \sparens{c}}{\deg(\lambda_i \genpoly_i)}$.
    Additionally, denote $\hat{\rankval} \definedas \relschmrank{\varpolyset}{\vec{\genpoly_h}} \leq \rankval$.
    It is enough to show that $\drelrank{\degree_M}{\variety}{\vec{\genpoly}} \leq 2 \hat{\rankval} + 2$,
    If $\deg(\vec{\genpoly}) < \degree_M$, then $\drelrank{\degree_M}{\variety}{\vec{\genpoly}} = 1 \leq 2 \rankval + 2$.
    Otherwise, if $\deg(\vec{\genpoly}) = \degree_M$, then the remark follows from Remark~\ref{remark-high-relative-rank-implies-high-relative-schmidt-rank}
    as:
    \[
        \drelrank{\degree_M}{\variety}{\vec{\genpoly}} =
        \relrank{\variety}{\vec{\genpoly}} \leq
        2 \cdot \relschmrank{\varpolyset}{\vec{\genpoly}} + 2
    \]
    Where:
    \[
        \relschmrank{\varpolyset}{\vec{\genpoly}}
        \relschmrank{\varpolyset}{\homopart{\vec{\genpoly}}} =
        \relschmrank{\varpolyset}{\homopart{\vec{\genpoly_h}}} =
        \relschmrank{\varpolyset}{\vec{\genpoly_h}} =
        \hat{\rankval}
    \]
\end{proof}


\newcommand{\etalchar}[1]{$^{#1}$}
\begin{thebibliography}{BFH{\etalchar{+}}13}

\bibitem[AGL24]{alrabiah2024randomlypuncturedreedsolomoncodes}
Omar Alrabiah, Venkatesan Guruswami, and Ray Li.
\newblock Randomly punctured reed--solomon codes achieve list-decoding capacity
  over linear-sized fields, 2024.

\bibitem[BDG24]{brakensiek2024generalizedgmmdspolynomialcodes}
Joshua Brakensiek, Manik Dhar, and Sivakanth Gopi.
\newblock Generalized gm-mds: Polynomial codes are higher order mds, 2024.

\bibitem[BDGZ24]{brakensiek2024agcodesachievelistdecoding}
Joshua Brakensiek, Manik Dhar, Sivakanth Gopi, and Zihan Zhang.
\newblock Ag codes achieve list-decoding capacity over constant-sized fields,
  2024.

\bibitem[BDM24]{yaacov2024sparsehighdimensionalexpanders}
Inbar {Ben Yaacov}, Yotam Dikstein, and Gal Maor.
\newblock Sparse high dimensional expanders via local lifts, 2024.

\bibitem[BFH{\etalchar{+}}13]{bhattacharyya2013locally}
Arnab Bhattacharyya, Eldar Fischer, Hamed Hatami, Pooya Hatami, and Shachar
  Lovett.
\newblock Every locally characterized affine-invariant property is testable,
  2013.

\bibitem[BGM24]{brakensiek2024genericreedsolomoncodesachieve}
Joshua Brakensiek, Sivakanth Gopi, and Visu Makam.
\newblock Generic reed-solomon codes achieve list-decoding capacity, 2024.

\bibitem[BHT13]{bhattacharyya2013algorithmic}
Arnab Bhattacharyya, Pooya Hatami, and Madhur Tulsiani.
\newblock Algorithmic regularity for polynomials and applications, 2013.

\bibitem[BL14]{bhowmick2014list}
Abhishek Bhowmick and Shachar Lovett.
\newblock List decoding reed-muller codes over small fields, 2014.

\bibitem[BL15]{DBLP:journals/corr/0001L15}
Abhishek Bhowmick and Shachar Lovett.
\newblock Bias vs structure of polynomials in large fields, and applications in
  effective algebraic geometry and coding theory.
\newblock {\em CoRR}, abs/1506.02047, 2015.

\bibitem[BM24]{bafna2024characterizingdirectproducttesting}
Mitali Bafna and Dor Minzer.
\newblock Characterizing direct product testing via coboundary expansion, 2024.

\bibitem[Bog05]{10.1145/1060590.1060594}
Andrej Bogdanov.
\newblock Pseudorandom generators for low degree polynomials.
\newblock In {\em Proceedings of the Thirty-Seventh Annual ACM Symposium on
  Theory of Computing}, STOC '05, page 21–30, New York, NY, USA, 2005.
  Association for Computing Machinery.

\bibitem[Bsh14]{10.1145/2554797.2554828}
Nader Bshouty.
\newblock Testers and their applications.
\newblock In {\em Proceedings of the 5th Conference on Innovations in
  Theoretical Computer Science}, ITCS '14, page 327–352, New York, NY, USA,
  2014. Association for Computing Machinery.

\bibitem[BV07]{4389478}
Andrej Bogdanov and Emanuele Viola.
\newblock Pseudorandom bits for polynomials.
\newblock In {\em 48th Annual IEEE Symposium on Foundations of Computer Science
  (FOCS'07)}, pages 41--51, 2007.

\bibitem[CTS13]{Cohen2013PseudorandomGF}
Gil Cohen and Amnon Ta-Shma.
\newblock Pseudorandom generators for low degree polynomials from algebraic
  geometry codes.
\newblock {\em Electron. Colloquium Comput. Complex.}, TR13, 2013.

\bibitem[DD24]{dikstein2024agreementtheoremshighdimensional}
Yotam Dikstein and Irit Dinur.
\newblock Agreement theorems for high dimensional expanders in the small
  soundness regime: the role of covers, 2024.

\bibitem[DGV24]{dwivedi2024optimalpseudorandomgeneratorslowdegree}
Ashish Dwivedi, Zeyu Guo, and Ben~Lee Volk.
\newblock Optimal pseudorandom generators for low-degree polynomials over
  moderately large fields, 2024.

\bibitem[Dik22]{dikstein2022newhighdimensionalexpanders}
Yotam Dikstein.
\newblock New high dimensional expanders from covers, 2022.

\bibitem[DS08]{4558818}
Zeev Dvir and Amir Shpilka.
\newblock Noisy interpolating sets for low degree polynomials.
\newblock In {\em 2008 23rd Annual IEEE Conference on Computational
  Complexity}, pages 140--148, 2008.

\bibitem[DTT22]{DBLP:journals/cc/DoronTT22}
Dean Doron, Amnon Ta{-}Shma, and Roei Tell.
\newblock On hitting-set generators for polynomials that vanish rarely.
\newblock {\em Comput. Complex.}, 31(2):16, 2022.

\bibitem[DV22]{derksen2022fooling}
Harm Derksen and Emanuele Viola.
\newblock Fooling polynomials using invariant theory.
\newblock In {\em 2022 IEEE 63rd Annual Symposium on Foundations of Computer
  Science (FOCS)}, pages 399--406. IEEE, 2022.

\bibitem[GJX17]{guruswami2017efficientlylistdecodablepuncturedreedmuller}
Venkatesan Guruswami, Lingfei Jin, and Chaoping Xing.
\newblock Efficiently list-decodable punctured reed-muller codes, 2017.

\bibitem[GK22a]{gotlib2022listagreementexpansioncoboundary}
Roy Gotlib and Tali Kaufman.
\newblock List agreement expansion from coboundary expansion, 2022.

\bibitem[GK22b]{gowers2022equidistributionhighrankpolynomialsvariables}
W.~T. Gowers and Thomas Karam.
\newblock Equidistribution of high-rank polynomials with variables restricted
  to subsets of $\mathbbm{F}_p$, 2022.

\bibitem[GKZ08]{10.1145/1374376.1374417}
Parikshit Gopalan, Adam~R. Klivans, and David Zuckerman.
\newblock List-decoding reed-muller codes over small fields.
\newblock In {\em Proceedings of the Fortieth Annual ACM Symposium on Theory of
  Computing}, STOC '08, page 265–274, New York, NY, USA, 2008. Association
  for Computing Machinery.

\bibitem[GT07a]{green2007distribution}
Ben Green and Terence Tao.
\newblock The distribution of polynomials over finite fields, with applications
  to the gowers norms, 2007.

\bibitem[GT07b]{green2007primescontainarbitrarilylong}
Ben Green and Terence Tao.
\newblock The primes contain arbitrarily long arithmetic progressions, 2007.

\bibitem[GX14]{6875485}
Venkatesan Guruswami and Chaoping Xing.
\newblock Hitting sets for low-degree polynomials with optimal density.
\newblock In {\em 2014 IEEE 29th Conference on Computational Complexity (CCC)},
  pages 161--168, June 2014.

\bibitem[HHL19]{book}
Hamed Hatami, Pooya Hatami, and Shachar Lovett.
\newblock {\em Higher-order Fourier Analysis and Applications}.
\newblock 01 2019.

\bibitem[HL11]{hatami2011higher}
Hamed Hatami and Shachar Lovett.
\newblock Higher-order fourier analysis of and the complexity of systems of
  linear forms.
\newblock {\em Geometric And Functional Analysis}, 21(6):1331--1357, 2011.

\bibitem[KL08]{kaufman2008worst}
Tali Kaufman and Shachar Lovett.
\newblock Worst case to average case reductions for polynomials, 2008.

\bibitem[KZ18]{kazhdan2018polynomial}
David Kazhdan and Tamar Ziegler.
\newblock Polynomial functions as splines, 2018.

\bibitem[KZ19]{kazhdan2019extendingweaklypolynomialfunctions}
David Kazhdan and Tamar Ziegler.
\newblock Extending weakly polynomial functions from high rank varieties, 2019.

\bibitem[KZ20]{kazhdan2020propertieshighranksubvarieties}
David Kazhdan and Tamar Ziegler.
\newblock Properties of high rank subvarieties of affine spaces, 2020.

\bibitem[Lov08]{10.1145/1374376.1374455}
Shachar Lovett.
\newblock Unconditional pseudorandom generators for low degree polynomials.
\newblock In {\em Proceedings of the Fortieth Annual ACM Symposium on Theory of
  Computing}, STOC '08, page 557–562, New York, NY, USA, 2008. Association
  for Computing Machinery.

\bibitem[Lov18]{DBLP:journals/corr/abs-1803-02523}
Shachar Lovett.
\newblock {MDS} matrices over small fields: {A} proof of the {GM-MDS}
  conjecture.
\newblock {\em CoRR}, abs/1803.02523, 2018.

\bibitem[Lu12]{6243404}
Chi-Jen Lu.
\newblock Hitting set generators for sparse polynomials over any finite fields.
\newblock In {\em 2012 IEEE 27th Conference on Computational Complexity}, pages
  280--286, 2012.

\bibitem[LZ21]{lampert2021relative}
Amichai Lampert and Tamar Ziegler.
\newblock Relative rank and regularization, 2021.

\bibitem[Sch80]{10.1145/322217.322225}
J.~T. Schwartz.
\newblock Fast probabilistic algorithms for verification of polynomial
  identities.
\newblock {\em J. ACM}, 27(4):701–717, oct 1980.

\bibitem[Sch85]{10.1007/BF02392473}
Wolfgang~M. Schmidt.
\newblock {The density of integer points on homogeneous varieties}.
\newblock {\em Acta Mathematica}, 154(3-4):243 -- 296, 1985.

\bibitem[TZ11]{tao2011inverse}
Terence Tao and Tamar Ziegler.
\newblock The inverse conjecture for the gowers norm over finite fields in low
  characteristic, 2011.

\bibitem[Vio08]{4558816}
Emanuele Viola.
\newblock The sum of d small-bias generators fools polynomials of degree d.
\newblock In {\em 2008 23rd Annual IEEE Conference on Computational
  Complexity}, pages 124--127, 2008.

\bibitem[YH18]{DBLP:journals/corr/abs-1803-03752}
Hikmet Yildiz and Babak Hassibi.
\newblock Optimum linear codes with support constraints over small fields.
\newblock {\em CoRR}, abs/1803.03752, 2018.

\bibitem[Zip79]{Zippel1979ProbabilisticAF}
Richard Zippel.
\newblock Probabilistic algorithms for sparse polynomials.
\newblock In {\em Symbolic and Algebraic Computation}, 1979.

\end{thebibliography}
\end{document}
